\documentclass[11pt]{article}

\usepackage[margin=1.1in]{geometry}

\usepackage{amssymb,mathtools,natbib}

% Anchored formal environments for causalsmith papers.
% First mandatory argument = obj_id (machine anchor joining the paper to the
% Lean crosswalk; invisible in the rendered PDF). Optional argument = name.
% Each environment also defines \label{obj:<obj_id>} so prose can use
% \cref{obj:T-1}; cleveref derives the printed kind and number from the target.
\usepackage{amsmath}
\usepackage{mathrsfs}
\usepackage{amsthm}
\usepackage{booktabs}
% Long displays may break across pages; let TeX stretch a little harder before an
% overfull \hbox so wide formulas/tables are less likely to run off the page.
\allowdisplaybreaks
\setlength{\emergencystretch}{3em}
\newtheorem{theorem}{Theorem}
\newtheorem{lemma}{Lemma}
\newtheorem{assumption}{Assumption}
\newtheorem{definition}{Definition}
\newtheorem{proposition}{Proposition}
\newtheorem{remark}{Remark}
\newtheorem{citedresult}{Cited result}
% The amsthm counter is `algorithmbox` (not `algorithm`) so a later `\usepackage{algorithm}` cannot clash.
\newcounter{algorithmbox}
\renewcommand{\thealgorithmbox}{\arabic{algorithmbox}}
\NewDocumentEnvironment{theoremv}{m o s}{\begin{theorem}\label{obj:#1}{\normalfont[\texttt{\detokenize{#1}}]\IfValueT{#2}{\ (#2)}\IfBooleanT{#3}{\verificationfootnotemark}\ }}{\end{theorem}}
\NewDocumentEnvironment{lemmav}{m o s}{\begin{lemma}\label{obj:#1}{\normalfont[\texttt{\detokenize{#1}}]\IfValueT{#2}{\ (#2)}\IfBooleanT{#3}{\verificationfootnotemark}\ }}{\end{lemma}}
\NewDocumentEnvironment{assumptionv}{m o s}{\begin{assumption}\label{obj:#1}{\normalfont[\texttt{\detokenize{#1}}]\IfValueT{#2}{\ (#2)}\IfBooleanT{#3}{\verificationfootnotemark}\ }}{\end{assumption}}
\NewDocumentEnvironment{definitionv}{m o s}{\begin{definition}\label{obj:#1}{\normalfont[\texttt{\detokenize{#1}}]\IfValueT{#2}{\ (#2)}\IfBooleanT{#3}{\verificationfootnotemark}\ }}{\end{definition}}
% A `propositionv` is a proved result tiered as a Proposition (theorem-grade, machine-verified).
\NewDocumentEnvironment{propositionv}{m o s}{\begin{proposition}\label{obj:#1}{\normalfont[\texttt{\detokenize{#1}}]\IfValueT{#2}{\ (#2)}\IfBooleanT{#3}{\verificationfootnotemark}\ }}{\end{proposition}}
% A `remarkv` holds NON-load-bearing interpretive / open-question content (e.g. a causalsmith `oeq:`
% open question). It is neither proved nor assumed; no theorem may depend on it.
\NewDocumentEnvironment{remarkv}{m o s}{\begin{remark}\label{obj:#1}{\normalfont[\texttt{\detokenize{#1}}]\IfValueT{#2}{\ (#2)}\IfBooleanT{#3}{\verificationfootnotemark}\ }}{\end{remark}}
% An `algorithmv` is a DEFINITION-kind object presented as a boxed, numbered-step procedure (an
% estimator / selection rule built in stages). Same anchor contract as `definitionv`; its body is
% ordinary LaTeX whose steps are `\begin{enumerate}` items, so tex2html needs no extra machinery.
% Presented in the CONVENTIONAL algorithm style readers expect from `algorithm`+`algorithmic`:
% a rule above, an "Algorithm N (Title)" caption line, a rule under the caption, the numbered
% steps, and a closing rule. It is NOT a float — the anchor contract requires the object to stay
% where the outline places it, and floats drift. The obj-id tag is suppressed here (it is
% bookkeeping, and a caption line carrying it reads as debug output in a finished paper).
\NewDocumentEnvironment{algorithmv}{m o s}%
  {\par\addvspace{\medskipamount}\noindent\rule{\linewidth}{0.8pt}\par\nobreak\vspace{2pt}%
   \refstepcounter{algorithmbox}\label{obj:#1}%
   \noindent\textbf{Algorithm \thealgorithmbox}\IfValueT{#2}{ \textnormal{(#2)}}%
   \IfBooleanT{#3}{\verificationfootnotemark}\par\nobreak\vspace{2pt}%
   \noindent\rule{\linewidth}{0.4pt}\par\nobreak\vspace{2pt}\normalfont}%
  {\par\nobreak\vspace{2pt}\noindent\rule{\linewidth}{0.8pt}\par\addvspace{\medskipamount}}
% A `citedv` is an IMPORTED external result the paper relies on but does NOT prove
% (a causalsmith `gate_class:"cited"` node). It is machine-checked against the cited
% source at F2.5, never proved here; the fixed provenance note makes that explicit
% so the environment can never be read as a contribution of this paper. The \citep
% attribution is supplied by the surrounding prose (P2), keyed to references.bib.
\NewDocumentEnvironment{citedv}{m o s}{\begin{citedresult}\label{obj:#1}{\normalfont[\texttt{\detokenize{#1}}]\IfValueT{#2}{\ (#2)}\IfBooleanT{#3}{\verificationfootnotemark}\ \textit{(imported result; machine-checked against the cited source, not proved here.)}\ }}{\end{citedresult}}
% \leanref{obj_id}{display text}: an inline first-mention link to a from-note object's
% verified Lean code. In the PDF it renders the display text plainly (the Lean link is a
% web-only affordance); tex2html turns it into a drawer-opening span.
\newcommand{\leanref}[2]{#2}
% For a theorem/lemma whose Lean declaration accepts a source-matched published proposition as an
% input, the starred anchored environment puts the mark in its heading; the generated text remains
% outside the frozen mathematical environment. Keep the combined macro for older paper bundles.
\newcommand{\verificationfootnotemark}{\footnotemark}
\newcommand{\verificationfootnotetext}[1]{\footnotetext{#1}}
\newcommand{\verificationfootnote}[1]{\verificationfootnotemark\verificationfootnotetext{#1}}
% xcolor before hyperref (hyperref would otherwise pull in plain `color` for colorlinks, and a
% later xcolor load clashes). The page-1 identity stamp at the end of this file needs a grey.
\usepackage{xcolor}
\usepackage{graphicx}
\usepackage{eso-pic}
% hyperref follows natbib and the environment macros; cleveref must follow hyperref.
\usepackage[colorlinks=true,linkcolor=blue,citecolor=blue,urlcolor=blue]{hyperref}
\usepackage[capitalise,nameinlink,noabbrev]{cleveref}
% Pin the names explicitly: labels are placed inside the underlying amsthm environments, and these
% declarations make the PDF contract independent of cleveref's theorem-name inference. House style:
% reference names always print with a capital first letter ("Theorem 1", "Definition 2"), in any
% sentence position — hence `capitalise` above and capitalized \crefname forms below.
\crefname{theorem}{Theorem}{Theorems}
\Crefname{theorem}{Theorem}{Theorems}
\crefname{lemma}{Lemma}{Lemmas}
\Crefname{lemma}{Lemma}{Lemmas}
\crefname{assumption}{Assumption}{Assumptions}
\Crefname{assumption}{Assumption}{Assumptions}
\crefname{definition}{Definition}{Definitions}
\Crefname{definition}{Definition}{Definitions}
\crefname{proposition}{Proposition}{Propositions}
\Crefname{proposition}{Proposition}{Propositions}
\crefname{remark}{Remark}{Remarks}
\Crefname{remark}{Remark}{Remarks}
\crefname{citedresult}{Cited result}{Cited results}
\Crefname{citedresult}{Cited result}{Cited results}
\crefname{algorithmbox}{Algorithm}{Algorithms}
\Crefname{algorithmbox}{Algorithm}{Algorithms}
% ---------------------------------------------------------------------------
% Page-1 identity stamp (arXiv style) and the pinned title-block date.
%
% Split of responsibility: THIS file owns the FORMAT (placement, font, colour, punctuation),
% the generated `paper_stamp.tex` owns only the VALUES (number+version, area, revised date,
% DOI, link target). P4 refreshes this template on every emit, so a layout fix reaches every
% bundle from a recompile; and because `paper_stamp.tex` is not one of the P2 assembly
% digest's authored inputs, a DOI that arrives after publication needs only that one small
% file rewritten plus a recompile — never a redraft, never a reassembly.
%
% Defaults first, so a bundle WITHOUT a stamp file compiles exactly as it always did:
% `\cspaperdate` falls back to `\today` and no stamp is drawn.
\providecommand*{\cspaperdate}{\today}  % the \date{} target
\providecommand*{\csstampid}{}          % e.g. CSWP-2026-008v3 (empty ⇒ no stamp at all)
\providecommand*{\csstamparea}{}        % e.g. Stat
\providecommand*{\csstampdate}{}        % e.g. 27 Aug 2026
\providecommand*{\csstampdoi}{}         % e.g. 10.5281/zenodo.123 (empty ⇒ DOI part omitted)
\providecommand*{\csstamplink}{}        % href target (DOI url, else the paper's page; may be empty)
\IfFileExists{paper_stamp.tex}{% GENERATED by CausalSmith P4 — paper identity stamp values. Do not hand-edit.
%
% Field order on the stamp (owned by paper_macros.tex):
%   <number>v<version> · doi:<doi> · [<Area>] <date>   — the DOI is never the last token, so a
%   text extractor cannot run it into whatever follows. Without a DOI that part is omitted.
% Format lives in paper_macros.tex; only the values are here, and this file is NOT one of
% the P2 assembly digest's authored inputs. A DOI arriving after publication therefore needs
% this file rewritten plus a `latexmk` recompile — never a redraft or a reassembly.
\renewcommand*{\csstampid}{CSWP-2026-023v5}
\renewcommand*{\csstamparea}{Partial ID}
\renewcommand*{\csstampdate}{29 Sep 2026}
\renewcommand*{\csstampdoi}{}
\renewcommand*{\csstamplink}{https://causalsmith.org/papers/pid_unlinked_propensity_ate_v1}
\renewcommand*{\cspaperdate}{29 September 2026}
}{}
\makeatletter
% Field order is deliberate: the DOI is NEVER the last token on the line. A trailing identifier
% runs into whatever a text extractor emits next, which makes it unsafe to match mechanically
% (the Zenodo stamp matcher reads exactly this line); the date is a safe terminator.
%   <number>v<version> · doi:<doi> · [<Area>] <date>      — with a DOI
%   <number>v<version> · [<Area>] <date>                  — without
\newcommand{\cs@stampsep}{\,\textperiodcentered\,}
\newcommand{\cs@stamptext}{%
  \csstampid
  \ifx\csstampdoi\@empty\else\cs@stampsep doi:\csstampdoi\fi
  \ifx\csstamparea\@empty\else\cs@stampsep[\csstamparea]\fi
  \ifx\csstampdate\@empty\else\quad\csstampdate\fi
}
\ifx\csstampid\@empty\else
  \definecolor{csstampgrey}{gray}{0.45}
  % Background shipout material: it occupies no space, so the text block and the \thanks
  % footnote are byte-identical to an unstamped compile. The starred form applies to the
  % NEXT page shipped only — i.e. page 1. Placed in the left margin (the text block starts
  % at 1.1in), reading bottom-to-top, in a Times-like face at 8pt.
  \AddToShipoutPictureBG*{%
    \AtPageLowerLeft{%
      \put(\LenToUnit{0.42in},\LenToUnit{1.1in}){%
        \rotatebox{90}{%
          \fontfamily{ptm}\fontsize{8}{10}\selectfont
          \ifx\csstamplink\@empty
            \textcolor{csstampgrey}{\cs@stamptext}%
          \else
            \expandafter\href\expandafter{\csstamplink}{\textcolor{csstampgrey}{\cs@stamptext}}%
          \fi
        }%
      }%
    }%
  }
\fi
\makeatother


\title{Assignment probabilities without row-level links: Identification, label design, and inference in randomized trials}

\author{CausalSmith\thanks{Machine-generated by the CausalSmith research pipeline. Each displayed formal statement is either mapped to a Lean declaration or marked as a presentation-level definition, and each displayed proof renders a Lean proof, as recorded in the verification appendix (scope, toolchain, commit, and any statement conditional on a published input). Cited work otherwise supplies framing and positioning.}}

\date{\cspaperdate}

\begin{document}

\maketitle

\begin{abstract}
We study a randomized trial in which treatment and outcome remain linked to a deterministic label of the assignment probability. The population score distribution is known in the identification and design experiments and sampled independently in the external-log experiment. Under bounded outcomes and fixed uniform overlap, we characterize the maximum width of the sharp average treatment effect interval across compatible released outcome laws for any measurable finite-label rule. Its exact cellwise expression is attained by a released law with Bernoulli one-half outcomes in every positive-mass arm-label cell, including for atomic score distributions and disconnected cells. Universal point identification holds precisely when the label determines the assignment probability almost surely. With a budget of \(K\) labels, optimal worst-case ambiguity has order \(K^{-1}\) over score laws with a common positive density lower bound; an equal-width construction gives the upper bound for every score law. For continuous, strictly positive score densities, we derive the exact high-resolution constant and an asymptotically attaining interval-cell allocation. For trial size \(n\) and a known score law, the minimax expected length of a uniformly honest interval has order \(K^{-1}+n^{-1/2}\) over the lower-density class. With a fixed release and an independent score log of size \(m\), a projection interval attains excess-length order \(n^{-1/2}+m^{-1/2}\); separate lower bounds establish both sampling directions at a fixed positive trial-to-log size ratio.
\end{abstract}

\section{Introduction}

Randomized assignment can use a probability that varies across people even when a released trial file records only treatment and outcome. We consider a release that attaches to each trial row a deterministic label of that assignment probability. In the identification, design, and known-score inference experiments, the analyst knows the population distribution of assignment probabilities. In the external-log experiment, that distribution is learned from an independent score sample. This observation scheme retains the link among label, treatment, and outcome, while the population score distribution supplies the probabilities represented within each label. The question is how much information the labels preserve about the average treatment effect (ATE), and how to choose them when their number is limited.

Under bounded potential outcomes and a fixed uniform overlap margin, the sharp ATE interval for a specified compatible released law follows from within-cell quantile rearrangement. We use the bounded-outcome specialization of \citet[Theorem 3.2 and eq. (3.1)]{FanShermanShum2014} for this fixed-law interval; \cref{obj:lem:fixed-released-law-baseline} records the resulting endpoints. Our release-design criterion then takes the largest interval width across the released outcome laws compatible with a given score distribution and label rule. \Cref{obj:thm:all-label-ambiguity} expresses that width as a sum of paired absolute-deviation losses over label cells and constructs two causal laws attaining its endpoints under the same released law. The construction applies to arbitrary score distributions on the overlap interval and measurable cells, including atoms and disconnected sets. \Cref{obj:thm:universal-point-identification} gives the corresponding disclosure criterion: worst-case ambiguity is zero exactly when the label determines the score almost surely.

This criterion turns finite-label release into a design problem measured in ATE units. For a budget of \(K\) labels, \cref{obj:thm:k-label-rate} bounds the smallest worst-case ambiguity above and below at order \(K^{-1}\) when the score density has a common positive lower bound. Its equal-width upper construction applies to every score distribution on the overlap interval. When the density is continuous and strictly positive, \cref{obj:thm:sharp-high-resolution-constant} gives the exact leading constant. Writing \(\varepsilon\) for the overlap margin, \(e\) for a score value, and \(f\) for the score density, an asymptotically attaining sequence places interval boundaries at equal increments of \(\int_\varepsilon^e \sqrt{f(t)/[t(1-t)]}\,dt\). Thus label allocation responds to both score mass and the sensitivity of inverse assignment weights.

The design criterion also governs inference from independent released trial rows. With a known score law and a jointly chosen release and uniformly honest interval, \cref{obj:thm:minimax-honest-length} establishes minimax expected length of order \(K^{-1}+n^{-1/2}\) over the common positive lower-density class, where \(n\) is trial size. An equal-width release and midpoint inverse-probability weights attain the upper order. A label budget proportional to \(\sqrt n\) balances ambiguity from score variation within cells with trial sampling variation. This connects the partial-identification criterion to confidence-set analysis \citep{ImbensManski2004,Stoye2009,ChernozhukovHongTamer2007,KaidoMolinariStoye2019}.

We also study a fixed finite-label release when the score distribution is learned from an independent log of \(m\) scores. The criterion here is worst-case expected positive excess length beyond the population sharp interval. For every measurable fixed release and score law on the overlap interval, \cref{obj:thm:external-score-root-order} constructs a uniformly honest projection interval with excess length bounded at order \(n^{-1/2}+m^{-1/2}\). A trial-sample lower bound holds at every positive pair of sample sizes, and a positive score-log lower bound holds asymptotically uniformly over trial sizes. Together they establish that order when the trial-to-log size ratio tends to a positive constant. The result includes atomic score laws and quantile functions with flat pieces.

The setting connects partial identification under combined marginals \citep{Manski1990,Tamer2010,Torgovitsky2019} with statistical matching \citep{DOrazio2006,Conti2017} and propensity-score subclassification \citep{Cochran1968,RosenbaumRubin1984}. The released label specifies a deterministic score cell, making within-cell association and the choice of cells central to the analysis. We document the scope of the Lean 4 verification in the appendix. We place the observation scheme among related work in the next section. The paper then develops the trial setup and sharp interval, characterizes information retained by labels, studies finite-label design and honest intervals, analyzes the independent score-log experiment, and closes with discussion and open questions.


\section{Related work}

Partial-identification analysis characterizes the target values compatible with an observed law \citep{Manski1990,HorowitzManski1995,Tamer2010}. For a known distribution of assignment probabilities and a fixed released trial law, the sharp average treatment effect interval developed below is the bounded-outcome, overlap specialization of \citet[Theorem 3.2 and eq. (3.1)]{FanShermanShum2014}. Within each released label, its endpoints account for the possible association between outcomes and assignment probabilities.

Data-combination and conditional-transport analyses also characterize outcomes consistent with observed marginal or conditional distributions. In particular, \citet[Section 3]{FanParkPassShi2025} obtain convex identified sets and transport-computable support functions for affine moment models with incomplete data. Here the fixed-law interval is one part of a release-design problem: the analysis ranges over compatible released trial laws to measure the largest attainable interval width, then chooses a finite-label rule to minimize that width. The score-based randomization model and this outer optimization are specific to the present problem.

Statistical matching and record-linkage studies analyze information recovered by linking separate files. \Citet[Section 4, especially Theorem 2]{Komarova2018} show that a positive bound on individual disclosure risk generally leaves their model parameter only partially pseudo-identified. Their object is a probabilistic matching rule based on personal identifiers and a disclosure constraint. Our released object is instead a deterministic function of the assignment score attached to every trial row; its information loss is measured by sharp ATE ambiguity, and point identification occurs exactly when the label determines the score almost surely.

Coarsening also appears in treatment-effect estimation with linked covariates. Under unconfoundedness, outcome smoothness, and restrictions on extreme scores, \citet[Assumptions 1--3 and Informal Theorem 1.1]{KalavasisMehrotraZampetakis2024} construct a data-dependent coarse IPW estimator whose error accounts for propensity-score inaccuracy and sampling. \Citet[Sections 3--5]{LeeWeidner2026} obtain treatment-effect bounds and inference by pooling outcomes across small groups of observations with linked covariates, including settings with limited overlap. For identification and known-score design, our experiment retains a deterministic label in place of the row-level score and assumes a known score marginal; the independent-score-log extension learns that marginal from a separate sample. Under a positive density floor, the resulting optimal ambiguity is of order \(K^{-1}\), and honest length in the known-score experiment is of order \(K^{-1}+n^{-1/2}\). These rates quantify label design under a different observation scheme and different regularity conditions.

The inference questions fit the literature on confidence sets for partially identified parameters \citep{ImbensManski2004,Stoye2009,ChernozhukovHongTamer2007,KaidoMolinariStoye2019}. With a known score distribution, the criterion is worst-case expected length for a jointly chosen label rule and uniformly honest interval. With an independent score log, the rule is fixed and the criterion is worst-case expected positive excess length beyond the population sharp interval from estimation with both released trial rows and the independent score log. These two criteria measure, respectively, the length of an honest interval under label coarsening and the additional length generated by the two sampled data sources.


\section{Trial setup and the sharp identified interval}

Throughout, generic positive constants may change from line to line, absolute values and Euclidean norms have their usual meanings, and \(O(\cdot)\) and \(o(\cdot)\) denote asymptotic bounds. Arm and label indices range over the sets defined below.

We consider \leanref{S-1}{a randomized causal population} governed by a law \leanref{sym:P}{\ensuremath{P}}. For an overlap margin \leanref{sym:epsilon}{\ensuremath{\varepsilon}}, the score \leanref{sym:E}{\ensuremath{E}} takes values in the interval \leanref{sym:Espace}{\ensuremath{\mathcal E}} given by \([\varepsilon,1-\varepsilon]\). The control and treatment potential outcomes \leanref{sym:Y0}{\ensuremath{Y_0}} and \leanref{sym:Y1}{\ensuremath{Y_1}} take values in the bounded outcome space \leanref{sym:Yspace}{\ensuremath{\mathcal Y=[0,1]}}. The assigned arm is \leanref{sym:A}{\ensuremath{A}}, and the observed outcome is \leanref{sym:Y}{\ensuremath{Y}}. The population score law is \leanref{sym:H}{\ensuremath{H}}; a measurable release rule \leanref{sym:g}{\ensuremath{g}} applied to the score determines the released label \leanref{sym:R}{\ensuremath{R}}; the joint law of the released trial variables is \leanref{sym:Prel}{\ensuremath{P_{\mathrm{rel}}}}. Our target is the average treatment effect \leanref{sym:theta}{\ensuremath{\theta(P)}} for the population law \(P\), defined by \(\theta(P)=\mathbb E_P[Y_1-Y_0]\). For identification, release design, and known-score inference, \(H\) is known to the analyst; in the external-log experiment it is learned from an independent sample of scores.




The population model combines overlap, a specified population score marginal, and score-based randomized assignment. The overlap margin is fixed throughout the analysis.

\begin{assumptionv}{ass:overlap}[Overlap parameter]
The overlap parameter \(\varepsilon\) satisfies \(0<\varepsilon<1/2\).
\end{assumptionv}

The uniform propensity-score overlap condition keeps both assignment probabilities positive throughout the score space \citep{FanShermanShum2014}. It also bounds the reciprocal assignment weights used in the endpoint formulas below.

\begin{assumptionv}{ass:score-marginal}[Score marginal law]
The score \(E\) has law \(H\) under \(P\):
\[
\mathcal L_P(E)=H.
\]
\end{assumptionv}

This specifies the population covariate or propensity marginal used in the data-combination experiment \citep{FanShermanShum2014}. It supplies the score distribution within each released label, even though individual scores are not linked to outcomes in the released data.

\begin{assumptionv}{ass:randomized-assignment}[Score-based treatment assignment]
\[
P(A=1\mid E,Y_0,Y_1)=E
\quad\text{almost surely}.
\]
\end{assumptionv}

Conditional on the score, assignment is Bernoulli and unconfounded with respect to both potential outcomes \citep{FanShermanShum2014}. Thus the score determines how the population score law is weighted in each treatment arm.

The observed outcome and the release are connected to the potential outcomes and score by two further conditions.

\begin{assumptionv}{ass:consistency}[Observed outcome consistency]
\[
Y=AY_1+(1-A)Y_0
\quad\text{almost surely}.
\]
\end{assumptionv}

Potential-outcome consistency identifies the observed outcome with the potential outcome under the assigned arm \citep{Rubin1974}. It connects each observed arm-label outcome distribution to the corresponding potential outcome.

\begin{assumptionv}{ass:deterministic-release}[Deterministic score release]
\[
R=g(E)
\quad\text{almost surely}.
\]
\end{assumptionv}

The deterministic release condition is specific to this analysis: each score is assigned a label by the declared rule. The released label therefore identifies a score cell while retaining uncertainty about the pairing of scores and outcomes within that cell.

\begin{assumptionv}{ass:released-law}[Released data law]
\[
\mathcal L_P(R,A,Y)=P_{\mathrm{rel}}.
\]
\end{assumptionv}

Observed-law compatibility requires a candidate causal law to reproduce the distribution available to the analyst, a standard restriction in data-combination identification \citep{FanShermanShum2014}. Here that distribution includes the label alongside assignment and outcome.

These conditions define the population laws under a release rule and, for a specified released distribution, the compatible laws.

\begin{definitionv}{def:causal-law-class}[Causal-law class \(\mathfrak P(H,g)\)]
The causal-law class \(\mathfrak P(H,g)\) consists of all laws \(P\) satisfying \cref{obj:ass:score-marginal,obj:ass:randomized-assignment,obj:ass:consistency,obj:ass:deterministic-release}.
\end{definitionv}

The class \leanref{sym:Pclass}{\ensuremath{\mathfrak P(H,g)}} in \cref{obj:def:causal-law-class} fixes the score law and release rule while allowing the joint distribution of scores and potential outcomes to vary subject to the stated population conditions.

\begin{definitionv}{def:compatible-law-class}[Compatible-law class \(\mathfrak M(H,g,P_{\mathrm{rel}})\)]
The compatible-law class is
\[
\mathfrak M(H,g,P_{\mathrm{rel}})
=
\bigl\{P\in\mathfrak P(H,g):\mathcal L_P(R,A,Y)=P_{\mathrm{rel}}\bigr\}.
\]
Here \(\mathfrak P(H,g)\) imposes \cref{obj:ass:score-marginal,obj:ass:randomized-assignment,obj:ass:consistency,obj:ass:deterministic-release}, and the released-law condition is \cref{obj:ass:released-law}.
\end{definitionv}

For a fixed released law, \leanref{sym:Mclass}{\ensuremath{\mathfrak M(H,g,P_{\mathrm{rel}})}} in \cref{obj:def:compatible-law-class} is the population class over which the identified set is formed.

We next fix the index sets used to describe released cells. Let \leanref{sym:J}{\ensuremath{J}} denote a positive label count; a later release-design problem uses a positive label budget \leanref{sym:K}{\ensuremath{K}}.

\begin{definitionv}{synth_3}[Released-label alphabets]
For each positive integer \(J\), the released-label alphabet is \(\mathcal R_J=\{1,\ldots,J\}\). In particular, \(\mathcal R_K=\{1,\ldots,K\}\) for a positive label budget \(K\); a release into \(\mathcal R_K\) may use fewer than \(K\) labels.
\end{definitionv}

The alphabet \leanref{sym:Rspace}{\ensuremath{\mathcal R_J}} in \cref{obj:synth_3} includes labels that may have zero probability under a particular release.

\begin{definitionv}{synth_2}[Treatment-arm set]
The treatment-arm set is \(\mathcal A=\{0,1\}\), where \(0\) denotes control and \(1\) denotes treatment.
\end{definitionv}

The arm index \leanref{sym:a}{\ensuremath{a}} ranges over \leanref{sym:Aspace}{\ensuremath{\mathcal A}} as defined in \cref{obj:synth_2}. For each arm-label cell, the observed outcome law and the score law induced by assignment enter the endpoint calculation together.
The endpoint formulas use the generalized quantile \(Q_F(u)=\inf\{x:F((-\infty,x])\ge u\}\) for \(0<u<1\), with support endpoints at \(u=0,1\). This convention also applies when \(F\) has atoms; its formal definition appears in \cref{obj:synth_4}.

\begin{definitionv}{def:arm-cell-primitives}[Arm-cell quantities $p_a,C_r,q_{ar},F_{ar},G_{ar},E_{ar},W_{ar}$]
For each arm \(a\in\mathcal A\) and label \(r\in\mathcal R_J\), define
\[
p_1(e)=e,\qquad p_0(e)=1-e,\qquad
C_r=\{e:g(e)=r\},\qquad
q_{ar}=\int_{C_r}p_a(e)\,H(de).
\]
When \(P_{\mathrm{rel}}(A=a,R=r)>0\), define the observed outcome law
\[
F_{ar}=\mathcal L_{P_{\mathrm{rel}}}(Y\mid A=a,R=r).
\]
For \(q_{ar}>0\), define
\[
G_{ar}(de)=\frac{\mathbf 1\{e\in C_r\}p_a(e)\,H(de)}{q_{ar}},
\qquad E_{ar}\sim G_{ar},
\qquad W_{ar}=\frac{1}{p_a(E_{ar})}.
\]
The endpoint formulas apply to triples \((H,g,P_{\mathrm{rel}})\) satisfying \(P_{\mathrm{rel}}(A=a,R=r)=q_{ar}\) for every arm and label. Their sums include cells with \(q_{ar}>0\), for which \(F_{ar}\) is defined.
\end{definitionv}

For each arm-label cell, the observed outcome law and the score law induced by assignment enter the endpoint calculation together.

\begin{definitionv}{def:mean-endpoints}[Mean endpoints $\underline\mu_a,\overline\mu_a$]
For a triple \((H,g,P_{\mathrm{rel}})\) satisfying \(P_{\mathrm{rel}}(A=a,R=r)=q_{ar}\) for every arm and label, define
\[
\underline\mu_a
=\sum_{r:q_{ar}>0}q_{ar}\int_0^1
Q_{F_{ar}}(u)Q_{W_{ar}}(1-u)\,du,
\qquad
\overline\mu_a
=\sum_{r:q_{ar}>0}q_{ar}\int_0^1
Q_{F_{ar}}(u)Q_{W_{ar}}(u)\,du.
\]
\end{definitionv}

In \cref{obj:def:arm-cell-primitives}, \leanref{sym:p}{\ensuremath{p_a(e)}} is the arm-specific assignment probability, \leanref{sym:C}{\ensuremath{C_r}} is the score cell assigned label \(r\), and \leanref{sym:q}{\ensuremath{q_{ar}}} is its arm-label probability. The observed outcome law \leanref{sym:F}{\ensuremath{F_{ar}}} and the assignment-weighted score law \leanref{sym:G}{\ensuremath{G_{ar}}} are the two marginals available for a positive-mass cell. The reciprocal weight \leanref{sym:W}{\ensuremath{W_{ar}}} converts an arm-specific outcome distribution into a potential-outcome mean; its association with outcomes within the cell determines the endpoint range.

\begin{definitionv}{def:sharp-ate-set}[Sharp ATE interval $I(P_{\mathrm{rel}};H,g)$]
\[
I(P_{\mathrm{rel}};H,g)
=
[\underline\mu_1-\overline\mu_0,\,
 \overline\mu_1-\underline\mu_0].
\]
\end{definitionv}

The generalized quantile \leanref{sym:Q}{\ensuremath{Q_F}} in \cref{obj:synth_4} makes the same endpoint expressions apply to discrete and continuous cell distributions. Pairing outcome and reciprocal-weight quantiles in opposite or matching order gives the lower and upper mean endpoints.

\begin{definitionv}{synth_4}[Generalized quantile]
For a probability law \(F\) with compact support in \(\mathbb R\), define
\[
Q_F(u)=\inf\{x\in\mathbb R:F((-\infty,x])\ge u\},\qquad 0<u<1.
\]
Set \(Q_F(0)=\inf\operatorname{supp}(F)\) and \(Q_F(1)=\sup\operatorname{supp}(F)\). At an atom \(x\), this convention gives \(Q_F(u)=x\) whenever \(F((-\infty,x))<u\le F((-\infty,x])\).
\end{definitionv}

The interval \leanref{sym:Isharp}{\ensuremath{I(P_{\mathrm{rel}};H,g)}} in \cref{obj:def:sharp-ate-set} combines the attainable arm-specific mean endpoints into the fixed-law sharp identified interval for the average treatment effect. The mathematical antecedent is \citet[Theorem 3.2 and eq. (3.1)]{FanShermanShum2014}; the displayed result here is proved under the paper's stated model conditions.

Before reporting that interval for a proposed release, one can check whether the declared score law and released law admit a common causal population. The observable criterion uses the label masses and the treated masses within labels.

\begin{theoremv}{prop:compatibility}[Compatibility of released laws]
Suppose \cref{obj:ass:overlap} holds. Let \(H\) be a probability law on the score interval \(\mathcal E=[\varepsilon,1-\varepsilon]\), let \(g:\mathcal E\to\mathcal R_J\) be a measurable release rule, and let \(P_{\mathrm{rel}}\) be a probability law on \((R,A,Y)\). The compatible causal-law class \(\mathfrak M(H,g,P_{\mathrm{rel}})\) of \cref{obj:def:compatible-law-class} is nonempty if and only if, for every \(r\in\mathcal R_J\),
\[
\begin{aligned}
P_{\mathrm{rel}}(R=r) &= H(C_r),\\
P_{\mathrm{rel}}(A=1,R=r) &= \int_{C_r} e\,H(de).
\end{aligned}
\]
\end{theoremv}

Under \cref{obj:ass:randomized-assignment,obj:ass:deterministic-release}, the score mass of a cell gives its label probability, and score-weighting that mass gives its treated probability. The control probability then follows from the label total. \Cref{obj:prop:compatibility} establishes that these equalities also suffice for a compatible causal law, so an archive can check the coherence of its declared release and randomization law before using the fixed-law interval.


\section{Information retained by a score label}

For a fixed score law and deterministic release, the sharp interval in \cref{obj:def:sharp-ate-set} can vary with the released outcome law. The worst-case ambiguity \leanref{sym:D}{\ensuremath{D_H(g)}}, the largest such interval width over the causal laws in \cref{obj:def:causal-law-class}, measures the information retained by the label across outcome laws.

Three criteria answer different questions. The sharp width measures identification for a specified released law; worst-case ambiguity maximizes that width over released laws at a fixed score law and release. With trial sampling and a known score law, \cref{obj:thm:minimax-honest-length} studies expected honest interval length. With an independently sampled score law, \cref{obj:def:external-excess-risk} studies expected length in excess of the population sharp interval. The latter two criteria include sampling uncertainty; the sharp width does not.

\begin{definitionv}{def:worst-case-ambiguity}[Worst-case ambiguity $D_H(g)$]
\[
D_H(g)
=
\sup_{P\in\mathfrak P(H,g)}
\operatorname{len}\!\left\{
I\bigl(\mathcal L_P(R,A,Y);H,g\bigr)
\right\}.
\]
\end{definitionv}

The supremum has a cellwise expression in terms of the score law and release rule. The next result also gives a released outcome law at which the width is attained.

\begin{theoremv}{thm:all-label-ambiguity}[Exact all-label ambiguity width]
Suppose that:
\begin{itemize}
\item \textbf{(Overlap.)} The overlap condition in \cref{obj:ass:overlap} holds.
\item \textbf{(Score law.)} \(H\) is a probability law on \(\mathcal E=[\varepsilon,1-\varepsilon]\).
\item \textbf{(Release.)} \(g:\mathcal E\to\mathcal R_J\) is measurable.
\end{itemize}
Write \(C_r=\{e:g(e)=r\}\), \(p_1(e)=e\), \(p_0(e)=1-e\), and \(q_{ar}=\int_{C_r}p_a(e)\,H(de)\). Let \(P_{\mathrm{rel}}\) be the released law
\[
P_{\mathrm{rel}}
=\sum_{r\in\mathcal R_J}\sum_{a\in\{0,1\}}
\frac{q_{ar}}{2}\bigl(\delta_{(r,a,0)}+\delta_{(r,a,1)}\bigr).
\]
Thus its observed outcome law \(F_{ar}\) is Bernoulli\((1/2)\) in every positive-mass arm-label cell.

The worst-case ambiguity \(D_H(g)\) of \cref{obj:def:worst-case-ambiguity} satisfies
\[
D_H(g)=
\sum_{r\in\mathcal R_J}
\left[
\inf_{c\in[1/(1-\varepsilon),\,1/\varepsilon]}
\int_{C_r}|1-ce|\,H(de)
+
\inf_{c\in[1/(1-\varepsilon),\,1/\varepsilon]}
\int_{C_r}|1-c(1-e)|\,H(de)
\right].
\]

For each \(a\in\{0,1\}\) and each cell with \(q_{ar}>0\), let \(E_{ar}\) have law \(p_a(e)\mathbf 1_{C_r}(e)H(de)/q_{ar}\), and let \(W_{ar}=1/p_a(E_{ar})\). Define
\[
\underline\mu_a
=\sum_{r:q_{ar}>0}q_{ar}\int_0^1
Q_{F_{ar}}(u)Q_{W_{ar}}(1-u)\,du,
\qquad
\overline\mu_a
=\sum_{r:q_{ar}>0}q_{ar}\int_0^1
Q_{F_{ar}}(u)Q_{W_{ar}}(u)\,du.
\]
The sharp ATE interval for \(P_{\mathrm{rel}}\) has length \(D_H(g)\). Moreover, there are two probability laws \(P_-,P_+\in\mathfrak M(H,g,P_{\mathrm{rel}})\), as defined in \cref{obj:def:compatible-law-class}, such that
\[
\mathbb E_{P_-}[Y_1-Y_0]=\underline\mu_1-\overline\mu_0,
\qquad
\mathbb E_{P_+}[Y_1-Y_0]=\overline\mu_1-\underline\mu_0,
\qquad
\mathbb E_{P_+}[Y_1-Y_0]-\mathbb E_{P_-}[Y_1-Y_0]=D_H(g).
\]
\end{theoremv}

Assignment tilts the score distribution within a label differently in the two arms. Pairing the Bernoulli outcome with reciprocal assignment weights in opposite or matching quantile order produces the mean endpoints in \cref{obj:thm:all-label-ambiguity}. Its two compatible causal laws attain those endpoints under the same released outcome law, so the paired absolute-deviation expression gives the exact worst-case width. Because the expression depends on \(H\) and \(g\), it permits comparison of candidate releases without specifying an outcome regression. The result applies to arbitrary probability laws \(H\), including atomic laws, and to measurable releases with disconnected score cells.

For a concrete one-label release, take \(\varepsilon=1/4\), let \(H\) put mass \(1/2\) on each of the scores \(1/4\) and \(3/4\), and set \(g\equiv1\). Use the released law in \cref{obj:thm:all-label-ambiguity}, so outcomes are Bernoulli\((1/2)\) in each arm. In either arm the lower and upper attainable potential-outcome means are \(1/3\) and \(2/3\): assignment tilts the two scores, but the outcome can be matched to the larger or smaller reciprocal weight. The sharp ATE interval is therefore \([-1/3,1/3]\), with width \(2/3\). If each score instead receives its own label, the score is recovered from the label and the ambiguity width is zero. This example isolates uncertainty from within-label score variation before any trial sampling is introduced.

The zero-width case has a direct disclosure interpretation. The following characterization expresses it both as recovery of the score from the label and as concentration of the score within every positive-mass cell.

\begin{theoremv}{thm:universal-point-identification}[Universal point identification]
Let \(J\in\mathbb N\), let \(H\) be a probability law on the score interval \(\mathcal E=[\varepsilon,1-\varepsilon]\), and let \(g:\mathcal E\to\mathcal R_J\) be a measurable deterministic release into \(J\) labels. Under \cref{obj:ass:overlap}, write \(C_r=\{e\in\mathcal E:g(e)=r\}\) for each label \(r\). Each of the following is equivalent to \(D_H(g)=0\), where \(D_H(g)\) is defined in \cref{obj:def:worst-case-ambiguity}:
\begin{enumerate}
\item There exists a measurable map \(h:\mathcal R_J\to\mathcal E\) such that \(e=h(g(e))\) for \(H\)-almost every \(e\).
\item For every \(r\in\mathcal R_J\) with \(H(C_r)>0\), there exists \(x\in\mathcal E\) such that \(e=x\) for \(H\)-almost every \(e\in C_r\).
\end{enumerate}
\end{theoremv}

The absolute-deviation terms in \cref{obj:thm:all-label-ambiguity} vanish when each positive-mass cell carries a single score value. Score variation within a cell produces positive worst-case ambiguity, while score-degenerate cells yield universal point identification. For a finite-support score law with at most \(J\) positive-mass atoms, giving each atom a distinct label achieves point identification. With more atoms than labels, at least one positive-mass cell contains score variation, so universal point identification fails.


\section{Designing a finite-label release}

The cellwise criterion in \cref{obj:thm:universal-point-identification} makes label design an allocation problem. A release with \leanref{sym:K}{\ensuremath{K}} labels groups score values, and \cref{obj:def:worst-case-ambiguity} measures the resulting worst-case ATE ambiguity. A lower bound on score density supports a comparison across score laws.

\begin{assumptionv}{ass:density-lower}[Lower bound on score density]
The score law \(H\) has density \(f\) on the score interval \(\mathcal E\), with
\[
H(de)=f(e)\,de,\qquad f(e)\ge m_f
\]
for Lebesgue-almost every \(e\in\mathcal E\).
\end{assumptionv}

The density \leanref{sym:f}{\ensuremath{f}} and its floor \leanref{sym:mf}{\ensuremath{m_f}}, when the floor is positive, give the standard compact-support positive-density condition used in quantization \citep{LemaireMontesPages2020}. A positive floor assigns every subinterval a baseline amount of score mass, which permits a uniform lower bound on ambiguity.

The release budget, its optimal ambiguity, and the common comparison class are defined as follows.

\begin{definitionv}{def:k-label-releases}[$K$-label release class $\mathcal G_K$]
\[
\mathcal G_K
=
\bigl\{g:\mathcal E\to\{1,\ldots,K\}:g\text{ is measurable}\bigr\}.
\]
\end{definitionv}

The release family \leanref{sym:GK}{\ensuremath{\mathcal G_K}} allows any measurable partition with at most \(K\) nonempty label cells. Its design criterion takes the infimum of worst-case ambiguity over that family.

\begin{definitionv}{def:optimal-k-ambiguity}[Optimal $K$-label ambiguity $\Delta_K(H)$]
\[
\Delta_K(H)=\inf_{g\in\mathcal G_K}D_H(g).
\]
\end{definitionv}

The optimal ambiguity \leanref{sym:DeltaK}{\ensuremath{\Delta_K(H)}} expresses the best worst-case ATE resolution available at the given label budget. For a uniform comparison, the density floor defines a class of score laws on the common overlap interval.

\begin{definitionv}{def:score-lower-density-class}[Lower-density score class $\mathcal H^{\mathrm{lb}}_{\varepsilon,m_f}$]
Under \cref{obj:ass:density-lower}, the class of score laws is
\[
\mathcal H^{\mathrm{lb}}_{\varepsilon,m_f}
=
\left\{
H\in\mathcal P(\mathcal E):
H(de)=f(e)\,de,\quad
f(e)\ge m_f\ \text{for Lebesgue-almost every }e\in\mathcal E
\right\}.
\]
\end{definitionv}

The class is nonempty only in the feasible range \(0<m_f\le (1-2\varepsilon)^{-1}\); this is the range in which the uniform lower bound below has content. Within this class, the label budget yields matching inverse-budget bounds. The constructive upper bound applies to every score law on the overlap interval.

\begin{theoremv}{thm:k-label-rate}[Bounds for optimal label ambiguity]
Suppose:
\begin{itemize}
\item \textbf{(Overlap.)} The condition in \cref{obj:ass:overlap} holds for \(\varepsilon\).
\item \textbf{(Density floor.)} \(m_f>0\).
\item \textbf{(Label budget.)} \(K\geq 1\) is an integer.
\end{itemize}
Let \(\ell=1-2\varepsilon\) and \(\mathcal E=[\varepsilon,1-\varepsilon]\). For every \(H\in\mathcal H^{\mathrm{lb}}_{\varepsilon,m_f}\) as defined in \cref{obj:def:score-lower-density-class},
\[
\frac{m_f\ell^2}{2(1-\varepsilon)K}\leq\Delta_K(H).
\]
For every probability law \(H\) on \(\mathcal E\), the equal-width release \(g:\mathcal E\to\mathcal R_K=\{1,\ldots,K\}\) defined by
\[
g(e)=1+\min\left\{\left\lfloor\frac{K(e-\varepsilon)}{\ell}\right\rfloor,K-1\right\}
\]
is measurable and satisfies
\[
\Delta_K(H)\leq\frac{\ell}{2\varepsilon(1-\varepsilon)K},
\qquad
D_H(g)\leq\frac{\ell}{2\varepsilon(1-\varepsilon)K}.
\]
\end{theoremv}

A positive density floor gives an inverse-budget lower bound on ambiguity for every \(K\)-label design. Equal-width cells give a matching upper bound for the optimal ambiguity under that floor, and an inverse-budget upper bound for any score law, including laws with atoms. The displayed constants thus translate a label budget into a worst-case ATE-width guarantee before outcomes are analyzed.

For a continuous, strictly positive score density, the leading coefficient has an exact law-specific form. An interval-cell release attains that coefficient asymptotically.

\begin{theoremv}{thm:sharp-high-resolution-constant}[Sharp high resolution constant]
Let \(\mathcal E=[\varepsilon,1-\varepsilon]\) be the score interval, and let \(H\) be a probability law on \(\mathcal E\). Suppose:
\begin{itemize}
\item \textbf{(Overlap.)} The overlap condition in \cref{obj:ass:overlap} holds.
\item \textbf{(Density.)} \(H(de)=f(e)\,de\) on \(\mathcal E\).
\item \textbf{(Regularity.)} The density \(f\) is continuous and strictly positive on \(\mathcal E\).
\end{itemize}
Let \(\Delta_K(H)\) and \(D_H(g)\) be the optimal \(K\)-label ambiguity and the worst-case ambiguity defined in \cref{obj:def:optimal-k-ambiguity,obj:def:worst-case-ambiguity}. For every positive integer \(K\), there is a measurable release \(g_K:\mathcal E\to\{0,\ldots,K-1\}\) whose label cells are intervals, such that
\[
\lim_{K\to\infty}K\Delta_K(H)
=
\lim_{K\to\infty}K D_H(g_K)
=
\frac14\left(
\int_{\varepsilon}^{1-\varepsilon}
\sqrt{\frac{f(e)}{e(1-e)}}\,de
\right)^2.
\]
\end{theoremv}

The theorem writes the attaining labels as \(0,\ldots,K-1\), whereas \(\mathcal G_K\) uses \(1,\ldots,K\). The bijection \(r\mapsto r+1\) preserves every label cell and hence preserves \(D_H(g_K)\), so the displayed construction is a member of the release class after relabeling.

The integrand weights score regions by their density and by the reciprocal-assignment factor \(1/[e(1-e)]\). It suggests an ordered interval allocation with boundaries at approximately equal increments of the cumulative weight \(\int_{\varepsilon}^{e}\sqrt{f(t)/[t(1-t)]}\,dt\): regions with greater weight receive narrower cells and more labels. The constant in \cref{obj:thm:sharp-high-resolution-constant} quantifies the leading ambiguity under optimal allocation, and its interval-cell sequence attains that value.


\section{Honest intervals with a known score law}

A finite-label release determines how much of the assignment score is retained. We now consider interval estimation when its score law is known to the analyst. For trial size \leanref{sym:n}{\ensuremath{n}}, the \leanref{S-2}{known-score trial experiment} observes independent released rows under a chosen release. The sampling condition is stated explicitly.

\begin{assumptionv}{ass:iid-trial-sampling}[Independent trial sampling]
The trial observations are independent and identically distributed under \(P\):
\[
(R_i,A_i,Y_i)_{i=1}^n\stackrel{\mathrm{iid}}{\sim}\mathcal L_P(R,A,Y).
\]
\end{assumptionv}

\begin{definitionv}{synth_1}[Average treatment effect functional]
For any measure \(P\) on full rows \((E,Y_0,Y_1,A,Y,R)\) with score parameter \(\varepsilon\) and \(J\) labels, define the average treatment effect functional \(\theta(P)\) by
\[
\theta(P)=\int (Y_1-Y_0)\,dP.
\]
\end{definitionv}

Independent superpopulation sampling governs how the released-row law generates the trial data and is a standard condition for interval inference in partially identified models \citep{ImbensManski2004}.

Coverage must account for every causal law in \cref{obj:def:causal-law-class} consistent with the chosen release and score law. At miscoverage level \leanref{sym:alpha}{\ensuremath{\alpha}}, the following criterion allows the release and interval procedure to be selected jointly.

In the next two displays, \(P^{\otimes n}\) acts only on functions of the observed released rows. Its pushforward to those rows is \(\mathcal L_P(R,A,Y)^{\otimes n}\), the sampling law in \cref{obj:ass:iid-trial-sampling}.

\begin{definitionv}{def:honest-procedures}[Uniformly honest procedures $\mathcal J_{n,K,\alpha}(H)$]
\[
\mathcal J_{n,K,\alpha}(H)
=
\left\{
(g,C_n):
g\in\mathcal G_K,\quad
\inf_{P\in\mathfrak P(H,g)}
P^{\otimes n}\!\left[
\theta(P)\in C_n\bigl((R_i,A_i,Y_i)_{i=1}^n\bigr)
\right]
\ge 1-\alpha
\right\}.
\]
\end{definitionv}

For each pair in \leanref{sym:Jhonest}{\ensuremath{\mathcal J_{n,K,\alpha}(H)}}, the interval \leanref{sym:Cn}{\ensuremath{C_n}} attains the stated coverage uniformly over causal laws for its release. Among these pairs, the design criterion is worst-case expected interval length.

\begin{definitionv}{def:minimax-honest-length}[Minimax honest length $\mathcal L^\star_{n,K,\alpha}(H)$]
\[
\mathcal L^\star_{n,K,\alpha}(H)
=
\inf_{(g,C_n)\in\mathcal J_{n,K,\alpha}(H)}
\sup_{P\in\mathfrak P(H,g)}
\mathbb E_{P^{\otimes n}}\!\left[\operatorname{len}(C_n)\right].
\]
\end{definitionv}

The criterion \leanref{sym:Lminimax}{\ensuremath{\mathcal L^\star_{n,K,\alpha}(H)}} measures the shortest worst-case expected length attainable through joint choice of a release and an honest interval. The next result gives matching orders over the positive lower-density class in \cref{obj:def:score-lower-density-class}, together with a sampling lower bound for every score law.

\begin{theoremv}{thm:minimax-honest-length}[Minimax honest interval length]
Consider the following conditions:
\begin{itemize}
\item \textbf{(Overlap.)} The overlap parameter \(\varepsilon\) satisfies \cref{obj:ass:overlap}.
\item \textbf{(Miscoverage.)} \(0<\alpha<1/2\).
\item \textbf{(Trial size.)} \(n\geq1\).
\item \textbf{(Label budget.)} \(K\geq1\).
\end{itemize}
There are constants \(C_{\varepsilon,\alpha}>0\) and \(c_{\varepsilon,\alpha}>0\), independent of \(H,n,K\), such that every score law \(H\) on \(\mathcal E=[\varepsilon,1-\varepsilon]\) satisfies
\[
c_{\varepsilon,\alpha}n^{-1/2}
\leq \mathcal L^\star_{n,K,\alpha}(H),
\]
where the minimax expected length is defined in \cref{obj:def:minimax-honest-length}.

For every \(m_f>0\), there is a constant \(c_{\varepsilon,m_f,\alpha}>0\) such that every \(H\in\mathcal H^{\mathrm{lb}}_{\varepsilon,m_f}\), as defined in \cref{obj:def:score-lower-density-class}, satisfies
\[
c_{\varepsilon,m_f,\alpha}\bigl(K^{-1}+n^{-1/2}\bigr)
\leq \mathcal L^\star_{n,K,\alpha}(H)
\leq C_{\varepsilon,\alpha}\bigl(K^{-1}+n^{-1/2}\bigr).
\]

An interval attaining the upper bound uses the equal-width release with labels \(r\in\mathcal R_K=\{1,\ldots,K\}\):
\[
g_K(e)=1+\min\left\{\left\lfloor\frac{K(e-\varepsilon)}{1-2\varepsilon}\right\rfloor,K-1\right\},
\qquad
z_r=\varepsilon+\frac{(r-\tfrac12)(1-2\varepsilon)}{K}.
\]
For released trial rows \((R_i,A_i,Y_i)_{i=1}^n\), define
\[
\widehat T_n=\frac1n\sum_{i=1}^n
\left\{\frac{A_iY_i}{z_{R_i}}-\frac{(1-A_i)Y_i}{1-z_{R_i}}\right\},
\qquad
R_{n,K}=\frac4\varepsilon\sqrt{\frac{\log(4/\alpha)}{n}}
+\frac{1-2\varepsilon}{2K\varepsilon(1-\varepsilon)}.
\]
With \(\operatorname{clampATE}(t)=\max\{-1,\min\{1,t\}\}\), the interval
\[
C_n(x)=
\left[
\operatorname{clampATE}\bigl(\widehat T_n(x)-R_{n,K}\bigr),
\operatorname{clampATE}\bigl(\widehat T_n(x)+R_{n,K}\bigr)
\right]
\]
satisfies \((g_K,C_n)\in\mathcal J_{n,K,\alpha}(H)\), with honesty defined in \cref{obj:def:honest-procedures}. For every \(P\in\mathfrak P(H,g_K)\), as defined in \cref{obj:def:causal-law-class}, its expected length under the \(n\)-fold released-row law satisfies
\[
\int |C_n(x)|\,dP_{\mathrm{rel}}^{\otimes n}(x)
\leq C_{\varepsilon,\alpha}\bigl(K^{-1}+n^{-1/2}\bigr).
\]
\end{theoremv}

The two terms in \cref{obj:thm:minimax-honest-length} reflect score variation within a label cell and sampling variation across released rows. Midpoint inverse-probability weights give an honest interval whose radius accounts for both. Over the positive lower-density class, a label budget of order \(\sqrt n\) balances these terms and attains root-\(n\) expected length. The matching lower bound makes this balance a property of the experiment, rather than of the particular interval construction.


\section{Inference with an independent score log}

The known-score analysis in \cref{obj:thm:minimax-honest-length} takes the score law as given. The \leanref{S-3}{external score-log experiment} instead fixes a finite-label release and combines the trial rows with an independent log of \leanref{sym:m}{\ensuremath{m}} scores. The score-log array \leanref{sym:Elog}{\ensuremath{(E_j^{\mathrm{log}})_{j=1}^m}} supplies information about the score law while retaining the trial’s released-row data structure. Two sampling conditions specify this combination.

\begin{assumptionv}{ass:external-log-iid}[Independent score-log sampling]
The external score-log observations are independent and identically distributed with law \(H\):
\[
(E_j^{\mathrm{log}})_{j=1}^m\stackrel{\mathrm{iid}}{\sim}H.
\]
\end{assumptionv}

Independent and identically distributed auxiliary-sample observations are a standard data-combination condition \citep{MarellaPfeffermann2019}. Here they make the empirical score distribution informative about the population score law.

\begin{assumptionv}{ass:external-log-independent}[Independent trial and score log]
The external score log is independent of the trial observations:
\[
(E_j^{\mathrm{log}})_{j=1}^m
\mathrel{\perp\!\!\!\perp}
(R_i,A_i,Y_i)_{i=1}^n.
\]
\end{assumptionv}

Independence across the two samples is the standard independent two-sample data-combination condition \citep{MarellaPfeffermann2019}. It gives the trial and score-log empirical measures separate sampling scales. The admissible law pairs and their joint sampling law can now be stated.

\begin{definitionv}{def:external-law-class}[External law class $\mathfrak X(g)$]
\[
\mathfrak X(g)
=
\left\{
(P,H):
H\in\mathcal P(\mathcal E),\quad
P\in\mathfrak P(H,g)
\right\}.
\]
The class uses the conditions in \cref{obj:ass:overlap,obj:ass:score-marginal,obj:ass:randomized-assignment,obj:ass:consistency,obj:ass:deterministic-release}.
\end{definitionv}

The external law class \leanref{sym:Xext}{\ensuremath{\mathfrak X(g)}} lets the score law vary over probability measures on the score interval, including atomic laws, while maintaining the causal and release conditions of \cref{obj:def:causal-law-class}.

\begin{definitionv}{def:external-experiment}[External experiment $\mathbb Q_{P,H}^{n,m}$]
For \((P,H)\in\mathfrak X(g)\), the joint law of \(n\) released trial observations \((R_i,A_i,Y_i)\) and \(m\) independent score-log observations \(E_j^{\mathrm{log}}\) is
\[
\mathbb Q_{P,H}^{n,m}
=
\mathcal L_P(R,A,Y)^{\otimes n}\otimes H^{\otimes m}.
\]
\end{definitionv}

The product law \leanref{sym:Qext}{\ensuremath{\mathbb Q_{P,H}^{n,m}}} makes coverage a joint-sample requirement. For a fixed release, honest intervals must cover the population average treatment effect uniformly over the external law class.

\begin{definitionv}{def:external-honest-procedures}[External honest procedures $\mathcal J^{\mathrm{ext}}_{n,m,\alpha}(g)$]
\[
\mathcal J^{\mathrm{ext}}_{n,m,\alpha}(g)
=
\left\{
C_{n,m}:
\inf_{(P,H)\in\mathfrak X(g)}
\mathbb Q_{P,H}^{n,m}\!\left[
\theta(P)\in C_{n,m}
\right]
\ge 1-\alpha
\right\}.
\]
\end{definitionv}

Membership in \leanref{sym:Jext}{\ensuremath{\mathcal J^{\mathrm{ext}}_{n,m,\alpha}(g)}} requires coverage for every admissible score law and causal law. Interval length is evaluated relative to the sharp identified interval in \cref{obj:def:sharp-ate-set} at the same released and score laws.

\begin{definitionv}{def:external-excess-risk}[Minimax external excess risk $\mathcal R^{\star}_{n,m,\alpha}(g)$]
\[
\mathcal R^{\star}_{n,m,\alpha}(g)
=
\inf_{C_{n,m}\in\mathcal J^{\mathrm{ext}}_{n,m,\alpha}(g)}
\sup_{(P,H)\in\mathfrak X(g)}
\mathbb E_{\mathbb Q_{P,H}^{n,m}}
\!\left[
\left(
\operatorname{len}(C_{n,m})
-
\operatorname{len}I(\mathcal L_P(R,A,Y);H,g)
\right)_+
\right].
\]
\end{definitionv}

The risk \leanref{sym:Rext}{\ensuremath{\mathcal R^{\star}_{n,m,\alpha}(g)}} measures expected positive excess length beyond the identified interval, then takes the worst case over admissible law pairs and the best honest procedure. To compare experiments with different trial and log sizes, the following normalization uses the sum of their square-root sampling scales.

\begin{definitionv}{def:external-normalized-risk}[Normalized external risk $T_{n,m,\alpha}(g)$]
\[
b_{n,m}=\left(n^{-1/2}+m^{-1/2}\right)^{-1},
\qquad
T_{n,m,\alpha}(g)=b_{n,m}\mathcal R^{\star}_{n,m,\alpha}(g).
\]
\end{definitionv}

The normalizer \leanref{sym:bext}{\ensuremath{b_{n,m}}} puts the excess risk on the two-sample root scale; \leanref{sym:Text}{\ensuremath{T_{n,m,\alpha}(g)}} records the resulting normalized risk. Constructing an interval on this scale uses distances between finite measures that account for both cell mass and distribution within a cell.

\begin{definitionv}{synth_5}[Finite-measure distances]
For finite nonnegative measures \(\mu,\nu\) on a compact interval \(S\), define their total-mass-plus-integrated-CDF distance by
\[
d_S(\mu,\nu)
=
|\mu(S)-\nu(S)|
+
\int_S
\bigl|\mu(S\cap(-\infty,t])-\nu(S\cap(-\infty,t])\bigr|\,dt.
\]
The outcome-measure distance \(d_{[0,1]}\) uses \(S=[0,1]\), with integration over \([0,1]\); the score-measure distance \(d_{\mathcal E}\) uses \(S=\mathcal E=[\varepsilon,1-\varepsilon]\), with integration over \([\varepsilon,1-\varepsilon]\).
\end{definitionv}

The total-mass term records arm-label frequencies, while the integrated distribution-function term records variation in outcome or score values. The projection below uses the trial and score-log empirical measures to select nearby candidate laws. Their compatible ATE ranges are collected into one confidence interval. The countable rational sieve makes this mathematical construction measurable.

\begin{definitionv}{def:external-projection-handle}[External projection confidence interval]
The external projection confidence interval \(C_{n,m}\) is defined for \(\varepsilon,\alpha\in\mathbb R\), \(J,n,m\in\mathbb N\), a release map \(g:[\varepsilon,1-\varepsilon]\to\mathcal R_J\), and a sample of \(n\) trial rows and \(m\) score-log values. Here \(\mathcal R_J\) is the set of \(J\) labels. For each arm \(a\) and label \(r\), form the empirical finite measures
\[
\widehat\lambda_{ar}(B)
=\frac1n\sum_{i=1}^{n}\mathbf 1\{R_i=r,A_i=a,Y_i\in B\},
\qquad
\widehat\sigma_{ar}(B)
=\frac1m\sum_{j=1}^{m}p_a(E_j^{\mathrm{log}})
\mathbf 1\{g(E_j^{\mathrm{log}})=r,E_j^{\mathrm{log}}\in B\}.
\]
Equip finite measures on each compact support with total-mass plus integrated-CDF distance, and set
\[
r_n=\frac{4J}{\alpha\sqrt n},
\qquad
r_m=\frac{2J(2-2\varepsilon)}{\alpha\sqrt m}.
\]

For \(q\in\mathbb Q\), define the clamped outcome and score support codes
\[
x(q)=\max\{0,\min\{1,q\}\},
\qquad
e_\varepsilon(q)=\varepsilon+(1-2\varepsilon)x(q).
\]
For a finite sequence \(L=((q_s,w_s))_{s=1}^{N}\) of rational support-weight pairs, define the finitely supported measures
\[
\nu_Y(L)=\sum_{s=1}^{N}\max\{w_s,0\}\,\delta_{x(q_s)},
\qquad
\nu_E(L)=\sum_{s=1}^{N}\max\{w_s,0\}\,\delta_{e_\varepsilon(q_s)},
\]
where the nonnegative weights are viewed as extended-real weights.

When \(0<\varepsilon<1/2\), let \(\mathcal D_{\varepsilon,J}\) consist of all pairs of arrays
\[
\left(
(\nu_Y(L^Y_{ar}))_{a,r},
(\nu_E(L^E_{ar}))_{a,r}
\right)
\]
formed from separate, arbitrary finite sequences \(L^Y_{ar}\) and \(L^E_{ar}\) for every arm-label cell, subject to equality of the outcome and score measures' total masses in each cell. For all other \(\varepsilon\), set \(\mathcal D_{\varepsilon,J}=\varnothing\). Define the candidate set by the two closed empirical-ball conditions
\[
\mathcal B_{n,m}
=
\left\{
(\lambda',\sigma')\in\mathcal D_{\varepsilon,J}:
\sum_{a,r}d_{[0,1]}(\lambda'_{ar},\widehat\lambda_{ar})\le r_n,\quad
\sum_{a,r}d_{\mathcal E}(\sigma'_{ar},\widehat\sigma_{ar})\le r_m
\right\},
\]
where \(\mathcal E=[\varepsilon,1-\varepsilon]\) is the score interval.

For any sieve array \(b=(\lambda^b,\sigma^b)\), put \(q^b_{ar}=\lambda^b_{ar}([0,1])=\sigma^b_{ar}(\mathcal E)\). In a cell with \(q^b_{ar}>0\), let \(F^b_{ar}=\lambda^b_{ar}/q^b_{ar}\) be the normalized outcome law and let \(W^b_{ar}\) be the pushforward of \(\sigma^b_{ar}/q^b_{ar}\) under \(e\mapsto 1/p_a(e)\). Define the lower and upper arm endpoints by
\[
\underline M_a(b)=\sum_{r:q^b_{ar}>0}q^b_{ar}\int_0^1 Q_{F^b_{ar}}(u)Q_{W^b_{ar}}(1-u)\,du,
\qquad
\overline M_a(b)=\sum_{r:q^b_{ar}>0}q^b_{ar}\int_0^1 Q_{F^b_{ar}}(u)Q_{W^b_{ar}}(u)\,du.
\]
Each zero-mass cell contributes \(0\) to both sums; no normalized law is needed for that cell. For each \(b\in\mathcal B_{n,m}\), set
\[
I(b)=[\underline M_1(b)-\overline M_0(b),\,\overline M_1(b)-\underline M_0(b)].
\]
When \(\mathcal B_{n,m}\ne\varnothing\), form
\[
J_{n,m}
=
[-1,1]\cap
\operatorname{cl}\operatorname{conv}
\left(\bigcup_{b\in\mathcal B_{n,m}}I(b)\right).
\]
The totalized interval is
\[
C_{n,m}
=
\begin{cases}
J_{n,m},&
\mathcal B_{n,m}\ne\varnothing\ \text{and}\ J_{n,m}\ne\varnothing,\\
\{0\},&
\mathcal B_{n,m}=\varnothing\ \text{or}\
\bigl(\mathcal B_{n,m}\ne\varnothing\ \text{and}\ J_{n,m}=\varnothing\bigr).
\end{cases}
\]
\end{definitionv}

In \cref{obj:def:external-projection-handle}, trial observations determine empirical outcome measures and the log determines assignment-weighted empirical score measures. The sieve combines finitely supported rationally coded measures with matched outcome and score masses in each arm-label cell, including arrays assembled independently across cells, and retains candidates in the two displayed empirical balls. This conservative projection contains the population sharp interval on the coverage event; closing and convexifying the candidate intervals produces the interval \leanref{sym:Cext}{\ensuremath{C_{n,m}}} used in the rate result.

\begin{theoremv}{thm:external-score-root-order}[External-score root order]
Let \(\mathcal E=[\varepsilon,1-\varepsilon]\), and let \(g:\mathcal E\to\mathcal R_J\) be a measurable finite-label release. Suppose:

\begin{itemize}
\item \textbf{(Overlap.)} The overlap condition in \cref{obj:ass:overlap} holds.
\item \textbf{(Miscoverage.)} \(0<\alpha<1/2\).
\item \textbf{(Joint sampling.)} For every \(n,m\), a joint sampling law \(\mathbb Q_{P,H}^{n,m}\) is specified for each admissible causal law \(P\) and score law \(H\). Its trial marginal consists of \(n\) independent draws from the released law. Its score-log marginal and independence from the trial sample satisfy \cref{obj:ass:external-log-iid,obj:ass:external-log-independent}.
\end{itemize}

Set
\[
L_\alpha=1+(1-2\alpha)^2,\qquad
S_g=
\sup_{\substack{r\in\mathcal R_J,\ e_1<e_2<e_3\\
e_1,e_2,e_3\in C_r}}
\frac{(e_3-e_2)(1-e_1/e_2)}
{\sqrt{(e_3-e_2)^2+(e_3-e_1)^2+(e_2-e_1)^2}},
\]
and
\[
c_{\mathrm{tr}}(\alpha)
=\frac{(1-2\alpha)\sqrt{\log L_\alpha}}{4},
\qquad
c_{\mathrm{log}}(g,\alpha)
=\frac{(1-2\alpha)\sqrt{\log L_\alpha}}{2\sqrt3}\,S_g.
\]
Then \(0<S_g\le 1\). There is a constant \(C_{\varepsilon,J,\alpha,g}>0\) such that, for every positive \(n,m\), the minimax excess risk of \cref{obj:def:external-excess-risk} satisfies
\[
c_{\mathrm{tr}}(\alpha)n^{-1/2}
\le \mathcal R^\star_{n,m,\alpha}(g)
\le C_{\varepsilon,J,\alpha,g}
\bigl(n^{-1/2}+m^{-1/2}\bigr).
\]
For each such \(n,m\), an interval procedure \(C_{n,m}\) belongs to the honest class of \cref{obj:def:external-honest-procedures}, agrees pointwise as a closed interval with the projection interval of \cref{obj:def:external-projection-handle}, and, for every admissible pair \((P,H)\), satisfies
\[
\mathbb E_{\mathbb Q_{P,H}^{n,m}}
\!\left[
\max\!\left\{0,\,
\operatorname{length}(C_{n,m})
-\operatorname{length}\bigl(I(P_{\mathrm{rel}};H,g)\bigr)
\right\}
\right]
\le C_{\varepsilon,J,\alpha,g}
\bigl(n^{-1/2}+m^{-1/2}\bigr).
\]

In the external-score direction,
\[
\liminf_{m\to\infty}\inf_{n\ge1}
\sqrt m\,\mathcal R^\star_{n,m,\alpha}(g)
\ge c_{\mathrm{log}}(g,\alpha)>0.
\]
For every \(\rho>0\), the normalized risk \(T_{n,m,\alpha}(g)\) of \cref{obj:def:external-normalized-risk} obeys
\[
0<
\max\left\{
\frac{c_{\mathrm{tr}}(\alpha)}{1+\sqrt\rho},
\frac{c_{\mathrm{log}}(g,\alpha)\sqrt\rho}{1+\sqrt\rho}
\right\}
\le
\liminf_{\substack{n,m\to\infty\\ n/m\to\rho}}
T_{n,m,\alpha}(g)
\le
\limsup_{\substack{n,m\to\infty\\ n/m\to\rho}}
T_{n,m,\alpha}(g)
\le C_{\varepsilon,J,\alpha,g}.
\]
\end{theoremv}

The upper bound in \cref{obj:thm:external-score-root-order} follows the two sources of empirical variation in the projection: released trial rows contribute the \(n^{-1/2}\) scale, and assignment-weighted score-log measures contribute the \(m^{-1/2}\) scale. The theorem also gives separate lower bounds. The trial-sample lower bound applies at every positive pair of sample sizes, and the score-log lower bound remains positive asymptotically uniformly over trial sizes. For a positive limiting trial-to-log ratio \leanref{sym:rho}{\ensuremath{\rho}}, these lower bounds and the projection interval establish the joint order \(n^{-1/2}+m^{-1/2}\). The law class includes atomic score distributions and quantile functions with flat pieces.


\section{Discussion and open questions}

A score label linked to each trial row preserves the association between a measurable set of assignment probabilities, treatment, and outcome. Given a known score law, the arm-label outcome and score distributions determine the sharp mean endpoints in \cref{obj:def:mean-endpoints} and hence the average treatment effect interval in \cref{obj:def:sharp-ate-set}. The worst-case width in \cref{obj:def:worst-case-ambiguity} measures the information retained by a release in units of the target parameter. \Cref{obj:thm:universal-point-identification} characterizes when that information suffices for point identification across compatible released laws.

For score laws on \([\varepsilon,1-\varepsilon]\) with a common density floor \(m_f\) satisfying \(0<m_f\le(1-2\varepsilon)^{-1}\), \cref{obj:thm:k-label-rate} gives an order-\(K^{-1}\) rate for optimal worst-case ambiguity. Under the score-law conditions of \cref{obj:thm:sharp-high-resolution-constant}, its high-resolution allocation rule assigns finer cells where local variation contributes more to the ambiguity criterion. The rule connects the placement of labels to the sharp interval that an analyst can recover from the released data. \Cref{obj:thm:minimax-honest-length} establishes a \(K^{-1}+n^{-1/2}\) minimax honest-length rate for known score laws with a positive density lower bound, attained up to constants by equal-width labels.

\subsection{Open questions}

The projection interval in \cref{obj:def:external-projection-handle} is defined through a countable rational sieve. A finite truncation and optimization rule, with an error bound that preserves coverage and controls excess length for atomic laws and arbitrary measurable cells, remains to be developed.

The two-source bounds in \cref{obj:thm:external-score-root-order} establish the root scale for excess interval length with an independent score log. A sharp limit would quantify the combined contribution of trial and score-log sampling when their sizes grow at a fixed positive ratio. The question is stated for the full external-law class and, where appropriate, for a specified regular subclass.

\begin{remarkv}{oeq:external-score-projection}[External-score sharp limits]
Fix a measurable release \(g\), \(0<\alpha<1/2\), and \(\rho\in(0,\infty)\). The sharp asymptotic analysis over \(\mathfrak X(g)\), with confidence procedures in \(\mathcal J^{\mathrm{ext}}_{n,m,\alpha}(g)\), concerns
\[
\liminf_{\substack{n,m\to\infty\\ n/m\to\rho}}T_{n,m,\alpha}(g)
\qquad\text{and}\qquad
\limsup_{\substack{n,m\to\infty\\ n/m\to\rho}}T_{n,m,\alpha}(g).
\]
It asks whether these limits coincide and, if so, for their common sharp constant. If a common sharp limit fails over the full class, the corresponding question concerns an explicitly defined regular subclass of \(\mathfrak X(g)\), with both uniform honesty and the supremum in expected excess length restricted to that subclass. The subclass analysis specifies its regularity conditions and sharp limit for each \(\rho\in(0,\infty)\), using the external projection construction to obtain an attaining procedure or a counterexample.
\end{remarkv}

For a regular subclass, the honesty requirement and worst-case excess-length criterion in \cref{obj:oeq:external-score-projection} apply to the same collection of laws. The projection construction in \cref{obj:def:external-projection-handle} supplies a concrete route for studying attainment, while the fixed sample-size ratio keeps both sources of sampling variation in the asymptotic comparison.


\appendix

\section*{Appendices}

\section{Proofs and verification scope}

\subsection{Fixed-law bounds and compatibility}

The fixed released-law calculation precedes the optimization over released laws. We attribute its mathematical antecedent to \citet[Theorem 3.2 and eq. (3.1)]{FanShermanShum2014}. The formal declaration and proof displayed here are independently checked under the paper's stated model conditions. The compatibility criterion in \cref{obj:prop:compatibility} identifies the released laws to which the calculation applies. The first auxiliary result records the arm-label masses implied by a score law and release rule.

\begin{lemmav}{lem:released-arm-cell-masses}[Released arm-cell masses]
Let \(H\) be a measure on the score interval \(\mathcal E\), let \(g:\mathcal E\to\mathcal R_J\) be a release rule, and let \(P\in\mathfrak P(H,g)\) as in \cref{obj:def:causal-law-class}. Write \(P_{\mathrm{rel}}=\mathcal L_P(R,A,Y)\) for the released law. Then \(H\) and \(P_{\mathrm{rel}}\) are probability laws, and, for every \(a\in\mathcal A\) and \(r\in\mathcal R_J\),
\[
P_{\mathrm{rel}}(A=a,R=r)=q_{ar},
\]
where \(q_{ar}\) is the arm-cell mass defined in \cref{obj:def:arm-cell-primitives}.
\end{lemmav}

\begin{proof}[Proof of \cref{obj:lem:released-arm-cell-masses}]
The law \(P\) is a probability law by \cref{obj:def:causal-law-class}. Its score marginal is \(H\) by \cref{obj:ass:score-marginal}, and its released-record marginal is \(P_{\mathrm{rel}}\). Both marginals are therefore probability laws. % lean: CausalSmith.PartialID.UnlinkedPropensityAte.compatibleCellMasses_of_randomizedCausalLaw

Fix \(r\in\mathcal R_J\). By \cref{obj:ass:deterministic-release}, \(R=g(E)\) almost surely. Conditioning on \((E,Y_0,Y_1)\) and applying \cref{obj:ass:randomized-assignment} gives
\[
\begin{aligned}
P_{\mathrm{rel}}(A=1,R=r)
&=P(A=1,g(E)=r)\\
&=\mathbb E_P\!\left[E\,\mathbf 1_{\{g(E)=r\}}\right]\\
&=\int_{C_r}e\,H(de).
\end{aligned}
\] % lean: CausalSmith.PartialID.UnlinkedPropensityAte.compatibleCellMasses_of_randomizedCausalLaw

For the control arm, partition the event \(E\in C_r\) by treatment assignment. The score marginal condition and the preceding identity yield
\[
\begin{aligned}
P_{\mathrm{rel}}(A=0,R=r)
&=P(E\in C_r)-P(A=1,E\in C_r)\\
&=H(C_r)-\int_{C_r}e\,H(de)\\
&=\int_{C_r}(1-e)\,H(de)
=q_{0r}.
\end{aligned}
\] % lean: CausalSmith.PartialID.UnlinkedPropensityAte.compatibleCellMasses_of_randomizedCausalLaw
The integral subtraction is valid because the score is bounded on \(\mathcal E\). Finally, \(p_1(e)=e\) in \cref{obj:def:arm-cell-primitives}, so the treated-arm identity is \(P_{\mathrm{rel}}(A=1,R=r)=q_{1r}\). These two identities cover every \(a\in\mathcal A\). % lean: CausalSmith.PartialID.UnlinkedPropensityAte.compatibleCellMasses_of_randomizedCausalLaw
\end{proof}

These mass equalities connect the causal-law class to the compatibility criterion in \cref{obj:prop:compatibility}. For each compatible released law, the fixed-law result gives sharp bounds through quantile rearrangement within positive-mass arm-label cells.

\begin{lemmav}{lem:fixed-released-law-baseline}[Sharp fixed released law bounds]
Let \(H\) be a measure on the score interval \(\mathcal E=[\varepsilon,1-\varepsilon]\), let \(g:\mathcal E\to\mathcal R_J\) be a release rule, and let \(P_{\mathrm{rel}}\) be a measure on the released label, treatment arm, and outcome. Suppose:

\begin{itemize}
\item \textbf{(Overlap.)} The score support satisfies \cref{obj:ass:overlap}.
\item \textbf{(Release.)} The rule \(g\) is measurable.
\item \textbf{(Compatibility.)} \(\mathfrak M(H,g,P_{\mathrm{rel}})\ne\varnothing\).
\end{itemize}

For each arm \(a\in\{0,1\}\) and label \(r\), let \(C_r=\{e:g(e)=r\}\), let \(p_a(e)\) be the probability of arm \(a\) at score \(e\), and set \(q_{ar}=\int_{C_r}p_a(e)\,H(de)\). In each cell with \(q_{ar}>0\), let \(F_{ar}\) be the outcome law conditional on arm \(a\) and label \(r\) under \(P_{\mathrm{rel}}\), let \(E_{ar}\) have law \(q_{ar}^{-1}\mathbf 1_{C_r}(e)p_a(e)\,H(de)\), and let \(W_{ar}=1/p_a(E_{ar})\). Writing \(Q\) for a generalized quantile, define
\[
\underline\mu_a
=\sum_{r:q_{ar}>0}q_{ar}\int_0^1 Q_{F_{ar}}(u)Q_{W_{ar}}(1-u)\,du,
\qquad
\overline\mu_a
=\sum_{r:q_{ar}>0}q_{ar}\int_0^1 Q_{F_{ar}}(u)Q_{W_{ar}}(u)\,du.
\]
Then, for every \(a\in\{0,1\}\),
\[
\left\{\mathbb E_P[Y_a]:P\in\mathfrak M(H,g,P_{\mathrm{rel}})\right\}
=[\underline\mu_a,\overline\mu_a].
\]
The attainable average treatment effects form the interval defined in \cref{obj:def:sharp-ate-set}:
\[
\left\{\mathbb E_P[Y_1-Y_0]:P\in\mathfrak M(H,g,P_{\mathrm{rel}})\right\}
=I(P_{\mathrm{rel}};H,g)
=[\underline\mu_1-\overline\mu_0,\,
  \overline\mu_1-\underline\mu_0].
\]
Moreover, there exist \(P^-,P^+\in\mathfrak M(H,g,P_{\mathrm{rel}})\) such that
\[
\mathbb E_{P^-}[Y_1-Y_0]=\underline\mu_1-\overline\mu_0,
\qquad
\mathbb E_{P^+}[Y_1-Y_0]=\overline\mu_1-\underline\mu_0.
\]
\end{lemmav}

\begin{proof}[Proof of \cref{obj:lem:fixed-released-law-baseline}]
\noindent\emph{1. Construct compatible endpoint laws.}
Compatibility implies that \(H\) and \(P_{\mathrm{rel}}\) are probability laws and, for every arm \(a\) and label \(r\),
\[
P_{\mathrm{rel}}(A=a,R=r)=q_{ar}.
\]
By \cref{obj:ass:overlap}, \(\varepsilon\le p_a(e)\le 1-\varepsilon\) on \(\mathcal E\). Thus
\[
\varepsilon H(C_r)\le q_{ar}\le(1-\varepsilon)H(C_r),
\]
so a zero-mass arm cell is also \(H\)-null. Define, on \(\mathcal E\),
\[
w_a(e)=\frac{1}{p_a(e)}.
\]
For each \(q_{ar}>0\), let \(U\) be uniform on \((0,1)\) and define two laws on outcome–weight pairs:
\[
\pi_{ar}^{-}
=\mathcal L\bigl(Q_{F_{ar}}(U),Q_{W_{ar}}(1-U)\bigr),
\qquad
\pi_{ar}^{+}
=\mathcal L\bigl(Q_{F_{ar}}(U),Q_{W_{ar}}(U)\bigr).
\]
Both have marginals \(F_{ar}\) and \(\mathcal L(W_{ar})\). They can be lifted to laws \(\gamma_{ar}^{s}\) on score–outcome pairs, for \(s\in\{-,+\}\), satisfying
\[
\mathcal L_{\gamma_{ar}^{s}}(E)=G_{ar},
\qquad
\mathcal L_{\gamma_{ar}^{s}}\bigl(Y,w_a(E)\bigr)=\pi_{ar}^{s}.
\]
Indeed, \(\mathcal L(W_{ar})\) is the common weight marginal of \(\pi_{ar}^{s}\) and \(\mathcal L(E_{ar},w_a(E_{ar}))\). Conditional on that weight, couple their outcome and score coordinates. Conditional laws exist on these standard Borel spaces.

For either sign, form the assigned and potential-outcome laws
\[
N_a^{s}=\sum_{r:q_{ar}>0}q_{ar}\gamma_{ar}^{s},
\qquad
M_a^{s}(de,dy)=w_a(e)\,N_a^{s}(de,dy).
\]
The score marginal of \(N_a^{s}\) is \(p_a(e)H(de)\): on \(C_r\), its score marginal is \(q_{ar}G_{ar}(de)=\mathbf1_{C_r}(e)p_a(e)H(de)\), and zero-mass cells are \(H\)-null. Consequently,
\[
M_a^{s}(de,[0,1])=H(de),
\qquad
p_a(e)M_a^{s}(de,dy)=N_a^{s}(de,dy).
\]
For signs \(s_0,s_1\in\{-,+\}\), glue \(M_0^{s_0}\) and \(M_1^{s_1}\) along their common score marginal \(H\), obtaining a law of \((E,Y_0,Y_1)\) with those two score–outcome marginals. Conditional on the resulting variables, draw \(A=1\) with probability \(E\), and set \(R=g(E)\) and \(Y=Y_A\). Denote the resulting full law by \(P^{s_0s_1}\).

This construction has score law \(H\), score-based assignment, consistency, and deterministic release. For \(q_{ar}>0\) and measurable \(B\subseteq[0,1]\), it also gives
\[
P^{s_0s_1}(R=r,A=a,Y\in B)
=\int_{C_r\times B}p_a(e)\,M_a^{s_a}(de,dy)
=q_{ar}F_{ar}(B)
=P_{\mathrm{rel}}(R=r,A=a,Y\in B).
\]
Both sides vanish when \(q_{ar}=0\). Hence every \(P^{s_0s_1}\) belongs to the class in \cref{obj:def:compatible-law-class}. % lean: CausalSmith.PartialID.UnlinkedPropensityAte.endpointFullLaw_compatibleCausalLaw

\noindent\emph{2. Compute their means and treatment effects.}
The two quantile couplings give, for every positive-mass cell,
\[
\int yw\,\pi_{ar}^{-}(dy,dw)
=\int_0^1 Q_{F_{ar}}(u)Q_{W_{ar}}(1-u)\,du,
\qquad
\int yw\,\pi_{ar}^{+}(dy,dw)
=\int_0^1 Q_{F_{ar}}(u)Q_{W_{ar}}(u)\,du.
\]
Since \(M_a^{s}=w_aN_a^{s}\), summing these identities over cells yields
\[
\mathbb E_{P^{s_0s_1}}[Y_a]
=\sum_{r:q_{ar}>0}q_{ar}\int yw\,\pi_{ar}^{s_a}(dy,dw)
=
\begin{cases}
\underline\mu_a,&s_a=-,\\
\overline\mu_a,&s_a=+.
\end{cases}
\]
Set
\[
P_{\mathrm{lo}}=P^{--},\qquad P_{\mathrm{hi}}=P^{++},
\qquad P^-=P^{+-},\qquad P^+=P^{-+},
\]
where the first sign selects the control-arm coupling. Bounded outcomes permit subtraction of their expectations, so
\begin{equation}\label{eq:fixed-law-ate-endpoints}
\mathbb E_{P^-}[Y_1-Y_0]=\underline\mu_1-\overline\mu_0,
\qquad
\mathbb E_{P^+}[Y_1-Y_0]=\overline\mu_1-\underline\mu_0.
\end{equation}
% lean: CausalSmith.PartialID.UnlinkedPropensityAte.endpointFullLaw_armMean

\noindent\emph{3. Bound every compatible arm mean.}
Fix \(P\in\mathfrak M(H,g,P_{\mathrm{rel}})\). For \(q_{ar}>0\), define its conditional outcome–weight coupling by
\[
\pi_{ar,P}
=\mathcal L_P\bigl(Y_a,w_a(E)\mid A=a,R=r\bigr).
\]
Consistency and the released-law condition give its outcome marginal \(F_{ar}\); score-based assignment and score marginality give its weight marginal \(\mathcal L(W_{ar})\). The same assignment identity, now weighted by \(w_a(E)\), gives
\[
\mathbb E_P[Y_a]
=\sum_{r:q_{ar}>0}q_{ar}\int yw\,\pi_{ar,P}(dy,dw).
\]
Cells with \(q_{ar}=0\) contribute zero.

For any coupling \(\pi\) of \(F_{ar}\) and \(\mathcal L(W_{ar})\), the layer-cake identity writes \(\int yw\,d\pi\) as the integral of joint upper-tail probabilities. For thresholds \(\tau_Y,\tau_W\ge0\), those probabilities lie between the intersection bounds
\[
\max\!\left\{P(Y>\tau_Y)+P(W>\tau_W)-1,0\right\}
\le P(Y>\tau_Y,W>\tau_W)
\le \min\!\left\{P(Y>\tau_Y),P(W>\tau_W)\right\}.
\]
The couplings \(\pi_{ar}^{-}\) and \(\pi_{ar}^{+}\) attain the respective bounds: their quantile upper-tail events are oppositely ordered and nested, respectively. Integrating the tail bounds therefore gives
\[
\int_0^1 Q_{F_{ar}}(u)Q_{W_{ar}}(1-u)\,du
\le \int yw\,\pi_{ar,P}(dy,dw)
\le \int_0^1 Q_{F_{ar}}(u)Q_{W_{ar}}(u)\,du.
\]
All integrals are finite because \(0\le Y_a\le1\) and \(w_a(E)\le\varepsilon^{-1}\). Multiplication by \(q_{ar}\) and summation show that
\begin{equation}\label{eq:fixed-law-arm-bound}
\underline\mu_a\le\mathbb E_P[Y_a]\le\overline\mu_a.
\end{equation}
% lean: CausalSmith.PartialID.UnlinkedPropensityAte.compatible_armMean_mem_endpoints

\noindent\emph{4. Fill each arm interval.}
For any \(x\in[\underline\mu_a,\overline\mu_a]\), choose \(\rho_a\in[0,1]\) with
\[
x=(1-\rho_a)\underline\mu_a+\rho_a\overline\mu_a,
\]
including \(\rho_a=0\) when the endpoints coincide, and define
\[
P_{a,x}=(1-\rho_a)P_{\mathrm{lo}}+\rho_aP_{\mathrm{hi}}.
\]
Probability normalization, the score and released marginals, and the assignment identity are preserved by this mixture; consistency and deterministic release hold under both component laws and hence under the mixture. Thus \(P_{a,x}\) is compatible. Linearity of integration gives
\[
\mathbb E_{P_{a,x}}[Y_a]
=(1-\rho_a)\mathbb E_{P_{\mathrm{lo}}}[Y_a]
+\rho_a\mathbb E_{P_{\mathrm{hi}}}[Y_a]=x.
\]
Together with \cref{eq:fixed-law-arm-bound}, this proves the claimed attainable interval for each arm. % lean: CausalSmith.PartialID.UnlinkedPropensityAte.compatibleCausalLaw_mixture

\noindent\emph{5. Fill the treatment-effect interval.}
For every compatible \(P\), the two arm bounds and
\(\mathbb E_P[Y_1-Y_0]=\mathbb E_P[Y_1]-\mathbb E_P[Y_0]\) give
\[
\underline\mu_1-\overline\mu_0
\le\mathbb E_P[Y_1-Y_0]
\le\overline\mu_1-\underline\mu_0.
\]
Conversely, for any \(x\) in this interval, choose \(\rho_{\mathrm{ATE}}\in[0,1]\) such that
\[
x=(1-\rho_{\mathrm{ATE}})(\underline\mu_1-\overline\mu_0)
+\rho_{\mathrm{ATE}}(\overline\mu_1-\underline\mu_0),
\]
and set
\[
P_x=(1-\rho_{\mathrm{ATE}})P^-+\rho_{\mathrm{ATE}}P^+.
\]
The same mixture argument makes \(P_x\) compatible, and linearity gives
\(\mathbb E_{P_x}[Y_1-Y_0]=x\). The interval is \(I(P_{\mathrm{rel}};H,g)\) by \cref{obj:def:sharp-ate-set}; the compatible laws \(P^-\) and \(P^+\) attain its two endpoints as computed in \cref{eq:fixed-law-ate-endpoints}. % lean: CausalSmith.PartialID.UnlinkedPropensityAte.fixedReleasedLaw_sharp
\end{proof}

The attainable endpoints make the interval sharp for each compatible released law. With the compatibility characterization in \cref{obj:prop:compatibility}, this fixed-law result supplies the interval whose width is optimized in \cref{obj:def:worst-case-ambiguity}.

\subsection{Ambiguity and finite-label design}

The all-label result in \cref{obj:thm:all-label-ambiguity} applies that outer optimization to the fixed-law interval, while \cref{obj:thm:universal-point-identification} characterizes universal point identification. For finite-label design, the proof first expresses the paired distortion as an optimization over arbitrary measurable partitions. The next lemma establishes the high-resolution limit over arbitrary measurable partitions and constructs an asymptotically attaining sequence of ordered interval cells.



\begin{lemmav}{lem:paired-high-resolution-quantization}[Paired high-resolution quantization]
Let \(\mathcal I=[a,b]\), let \(S\) be an integer, and let \(\beta_s:\mathcal I\times\mathcal I\to\mathbb R\), \(s=1,\ldots,S\), satisfy:
\begin{itemize}
\item \textbf{(Interval.)} \(a<b\).
\item \textbf{(Number of components.)} \(S\geq1\).
\item \textbf{(Continuity.)} Each \(\beta_s\) is continuous on \(\mathcal I\times\mathcal I\).
\item \textbf{(Positivity.)} \(\beta_s(x,z)>0\) for every \(s\) and every \(x,z\in\mathcal I\).
\end{itemize}
For each positive integer \(k\), define
\[
V_k=
\inf_{\substack{(B_1,\ldots,B_k)\ \mathrm{measurable,\ pairwise\ disjoint}\\
\bigcup_{j=1}^{k}B_j=\mathcal I,\quad z_{js}\in\mathcal I}}
\sum_{j=1}^{k}\sum_{s=1}^{S}
\int_{B_j}\beta_s(x,z_{js})|x-z_{js}|\,dx,
\qquad
w(x)=\sum_{s=1}^{S}\beta_s(x,x).
\]
Then
\[
\lim_{k\to\infty}kV_k
=\frac14\left(\int_a^b\sqrt{w(x)}\,dx\right)^2.
\]
For each positive integer \(k\), let \(a=t_0<\cdots<t_k=b\) be boundaries satisfying
\[
\int_a^{t_j}\sqrt{w(x)}\,dx
=\frac{j}{k}\int_a^b\sqrt{w(x)}\,dx
\qquad (j=0,\ldots,k).
\]
Set \(B_j=[t_{j-1},t_j)\) for \(j<k\), \(B_k=[t_{k-1},t_k]\), and \(z_{js}=(t_{j-1}+t_j)/2\) for every \(s\). These cells form a measurable, pairwise-disjoint partition of \(\mathcal I\), every \(z_{js}\) belongs to \(\mathcal I\), and
\[
\lim_{k\to\infty}
k\sum_{j=1}^{k}\sum_{s=1}^{S}
\int_{B_j}\beta_s(x,z_{js})|x-z_{js}|\,dx
=\frac14\left(\int_a^b\sqrt{w(x)}\,dx\right)^2.
\]
\end{lemmav}

\begin{proof}[Proof of \cref{obj:lem:paired-high-resolution-quantization}]
Write
\[
f(x)=\sqrt{w(x)},\qquad \Lambda=\int_a^b f(x)\,dx.
\]
Continuity, positivity, and compactness give
\[
\underline\beta
=\min_{\substack{1\le s\le S\\x,z\in[a,b]}}\beta_s(x,z)>0,
\qquad
\underline f=\min_{x\in[a,b]}f(x)>0,
\qquad
\overline f=\max_{x\in[a,b]}f(x)<\infty.
\]

For the stated construction, the continuous function \(x\mapsto\int_a^x f(u)\,du\) is strictly increasing from \(0\) to \(\Lambda\). Equivalently, its boundaries can be written
\[
t_j=\inf\left\{x\in[a,b]:
\frac{j\Lambda}{k}\le\int_a^x f(u)\,du\right\},
\qquad j=0,\ldots,k.
\]
Thus \(t_0=a\), \(t_k=b\), the boundaries are ordered, and each cell has \(f\)-mass \(\Lambda/k\). Put \(h_j=t_j-t_{j-1}\) and \(u_j=(t_{j-1}+t_j)/2\). Then
\[
0<h_j\le\frac{\Lambda}{k\underline f},
\qquad
k\sum_{j=1}^k h_j^2
\le\frac{\Lambda}{\underline f}\sum_{j=1}^k h_j
=\frac{\Lambda(b-a)}{\underline f}.
\]
Uniform continuity of \(f\) and the \(\beta_s\), together with
\(\int_{B_j}|x-u_j|\,dx=h_j^2/4\), shows that, for every \(\varepsilon>0\) and all sufficiently large \(k\), uniformly over the cells,
\[
\left|w(u_j)h_j^2-\left(\frac{\Lambda}{k}\right)^2\right|
\le\varepsilon h_j^2,
\qquad
\left|
\sum_{s=1}^S\int_{B_j}\beta_s(x,u_j)|x-u_j|\,dx
-\frac{w(u_j)h_j^2}{4}
\right|
\le\varepsilon h_j^2.
\]
The bounded scaled sum of squared cell lengths therefore gives
\[
k\sum_{j=1}^k\sum_{s=1}^S
\int_{B_j}\beta_s(x,u_j)|x-u_j|\,dx
\longrightarrow\frac{\Lambda^2}{4}.
\]
% lean: Causalean.Mathlib.Analysis.Quantization.companding_scaled_cost_tendsto

For the lower bound, take any admissible measurable partition and reproduction points, and define its cost by
\[
\mathcal D_k(B,z)
=\sum_{j=1}^k\sum_{s=1}^S
\int_{B_j}\beta_s(x,z_{js})|x-z_{js}|\,dx.
\]
Fix \(0<\eta<1\). Choose \(e>0\) sufficiently small that
\[
e\le\min\{1,\underline\beta/2\},
\qquad
e(2\overline f+1+S)\le\eta S\underline\beta.
\]
Uniform continuity supplies \(\delta>0\) such that, for each cell \(j\), whenever \(x,y\in[a,b]\) satisfy \(|x-z_{js}|<\delta\) and \(|y-z_{js}|<\delta\) for every \(s=1,\ldots,S\), one has
\[
\bigl|\beta_s(x,z_{js})-\beta_s(y,y)\bigr|<e
\quad(1\le s\le S),
\qquad |f(x)-f(y)|<e.
\]
Set
\[
G_j=\{x\in B_j:|x-z_{js}|<\delta\ \text{for every }s\},
\qquad D_j=B_j\setminus G_j.
\]
For a nonempty \(G_j\), choose \(y_j\in G_j\) and set
\[
\gamma_{js}=\beta_s(y_j,y_j)-e,\qquad
Q_j=f(y_j)+e,\qquad
\Gamma_j=\sum_{s=1}^S\gamma_{js}.
\]
These coefficients are positive. On \(G_j\), \(\beta_s(x,z_{js})\ge\gamma_{js}\) and \(f(x)\le Q_j\). Moreover, the choice of \(e\), \(w(y_j)\ge S\underline\beta\), and \(f(y_j)\le\overline f\) give
\[
(1-\eta)Q_j^2
\le w(y_j)-Se
=\Gamma_j.
\]

For any measurable \(G\subset\mathbb R\) of finite measure and any \(v\in\mathbb R\), filling the points nearest \(v\) first yields
\[
\int_G|x-v|\,dx
\ge\int_0^{|G|/2}(|G|-2r)\,dr
=\frac{|G|^2}{4}.
\]
With \(v=\Gamma_j^{-1}\sum_s\gamma_{js}z_{js}\), the triangle inequality gives
\(\sum_s\gamma_{js}|x-z_{js}|\ge\Gamma_j|x-v|\).
Consequently, including the trivial case \(G_j=\varnothing\),
\[
\sum_{s=1}^S\int_{B_j}\beta_s(x,z_{js})|x-z_{js}|\,dx
\ge\frac{\Gamma_j|G_j|^2}{4}
\ge\frac{1-\eta}{4}\left(\int_{G_j}f(x)\,dx\right)^2.
\]
Summing and applying Cauchy--Schwarz, with disjointness and coverage of the cells, gives
\[
k\mathcal D_k(B,z)
\ge\frac{1-\eta}{4}
\left(\Lambda-\sum_{j=1}^k\int_{D_j}f(x)\,dx\right)^2.
\]
% lean: Causalean.Mathlib.Analysis.Quantization.partition_good_mass_lower

At every \(x\in D_j\), some reproduction point is at distance at least \(\delta\). Hence
\[
\underline\beta\delta\sum_{j=1}^k|D_j|
\le\mathcal D_k(B,z),
\qquad
\sum_{j=1}^k\int_{D_j}f(x)\,dx
\le\frac{\overline f}{\underline\beta\delta}\mathcal D_k(B,z).
\]
If \(k\mathcal D_k(B,z)\le\Lambda^2/4\), the omitted \(f\)-mass is therefore at most
\(\overline f\Lambda^2/(4\underline\beta\delta k)\); for larger costs the desired asymptotic lower bound is immediate. It follows uniformly over admissible partitions that
\[
\liminf_{k\to\infty}kV_k\ge\frac{1-\eta}{4}\Lambda^2.
\]
Letting \(\eta\downarrow0\) proves the lower bound. The constructed cost bounds \(V_k\) from above and has limit \(\Lambda^2/4\), so \(kV_k\) has that limit as well.
% lean: Causalean.Mathlib.Analysis.Quantization.optimal_scaled_cost_tendsto

Finally, the ordered boundaries with \(t_0=a\) and \(t_k=b\) make the stated half-open cells, with the last cell closed on the right, a measurable pairwise-disjoint partition of \([a,b]\). Each \(u_j\) lies between its two boundaries, hence in \([a,b]\), and serves as \(z_{js}\) for every component.
% lean: CausalSmith.PartialID.UnlinkedPropensityAte.pairedQuantization_tendsto
\end{proof}

The auxiliary result identifies the leading constant for the paired loss and gives an attaining sequence of ordered cells. It supplies the quantization component of the high-resolution conclusion in \cref{obj:thm:sharp-high-resolution-constant}.

\subsection{External-log excess risk}

The external-log analysis has two sampling contributions: released trial rows and independent score draws. The following certificate supplies lower bounds for both contributions to the excess-risk criterion in \cref{obj:def:external-excess-risk}.

\begin{lemmav}{lem:external-two-source-lower-certificate}[Two-source excess-risk lower bounds]
Let \(\mathcal E=[\varepsilon,1-\varepsilon]\) be the score interval, let \(\mathcal R_J\) be the finite label set, and fix a measurable release \(g:\mathcal E\to\mathcal R_J\). Let \(\mathbb Q_{P,H}^{n,m}\) be a family of joint sampling laws on \(n\) released trial rows and \(m\) score-log draws. Suppose:

\begin{itemize}
\item \textbf{(Overlap and level.)} \cref{obj:ass:overlap} holds and \(0<\alpha<1/2\).
\item \textbf{(Trial sampling.)} For every \(n,m\) and every admissible causal-law and score-law pair \((P,H)\) under \(g\), the trial marginal of \(\mathbb Q_{P,H}^{n,m}\) is the \(n\)-fold product of the released-row law under \(P\).
\item \textbf{(Score-log sampling.)} The joint sampling laws satisfy \cref{obj:ass:external-log-iid,obj:ass:external-log-independent} for every \(n,m\).
\end{itemize}

Set
\[
L_\alpha=1+(1-2\alpha)^2,\qquad
S_g=\sup_{\substack{r\in\mathcal R_J,\ e_1,e_2,e_3\in\mathcal E\\
e_1<e_2<e_3,\ g(e_1)=g(e_2)=g(e_3)=r}}
\frac{(e_3-e_2)(1-e_1/e_2)}
{\sqrt{(e_3-e_2)^2+(e_3-e_1)^2+(e_2-e_1)^2}},
\]
and
\[
c_{\mathrm{tr}}(\alpha)
=\frac{(1-2\alpha)\sqrt{\log L_\alpha}}{4},
\qquad
c_{\mathrm{log}}(g,\alpha)
=\frac{(1-2\alpha)\sqrt{\log L_\alpha}}{2\sqrt{3}}S_g.
\]
Then \(0<S_g\leq 1\). The external-score minimax excess risk of \cref{obj:def:external-excess-risk} satisfies, for every \(n,m\geq 1\),
\[
\mathcal R^\star_{n,m,\alpha}(g)
\geq \frac{c_{\mathrm{tr}}(\alpha)}{\sqrt n}.
\]
It also satisfies
\[
\liminf_{m\to\infty}\inf_{n\geq 1}
\left\{\sqrt m\,\mathcal R^\star_{n,m,\alpha}(g)\right\}
\geq c_{\mathrm{log}}(g,\alpha)>0.
\]
For every \(\rho>0\), the normalized risk of \cref{obj:def:external-normalized-risk} satisfies
\[
\liminf_{\substack{n,m\to\infty\\ n/m\to\rho}}
T_{n,m,\alpha}(g)
\geq
\max\left\{
\frac{c_{\mathrm{tr}}(\alpha)}{1+\sqrt\rho},
\frac{c_{\mathrm{log}}(g,\alpha)\sqrt\rho}{1+\sqrt\rho}
\right\}>0.
\]
\end{lemmav}

\begin{proof}[Proof of \cref{obj:lem:external-two-source-lower-certificate}]
Write \(\beta=1-2\alpha\), so \(0<\beta<1\), \(L_\alpha=1+\beta^2>1\), and both displayed certificate coefficients are positive once \(S_g>0\).

The interval \(\mathcal E\) is infinite, whereas the label set is finite. Thus some fiber of \(g\) contains three points \(e_1<e_2<e_3\). For any such triple, put
\[
D(e_1,e_2,e_3)
=\sqrt{(e_3-e_2)^2+(e_3-e_1)^2+(e_2-e_1)^2}.
\]
Overlap gives \(0<e_1/e_2<1\), and \(D(e_1,e_2,e_3)\ge e_3-e_2>0\). Consequently every ratio in the definition of \(S_g\) is positive and at most one. The set of ratios is nonempty, so
\[
0<S_g\le1,\qquad c_{\mathrm{tr}}(\alpha)>0,\qquad
c_{\mathrm{log}}(g,\alpha)>0.
\]
% lean: CausalSmith.PartialID.UnlinkedPropensityAte.fiberTripleSeparation_bounds

We first establish a uniform upper bound used to pass the score-log certificate through an infimum and a liminf. For finite measures on \(\mathcal V=[v_-,v_+]\), define
\[
d_{\mathcal V}(\mu,\nu)
=|\mu(\mathcal V)-\nu(\mathcal V)|
 +\int_{v_-}^{v_+}
   |\mu(\mathcal V\cap(-\infty,t])-\nu(\mathcal V\cap(-\infty,t])|\,dt.
\]
Let \(\lambda_{ar}(B)=P_{\mathrm{rel}}(A=a,R=r,Y\in B)\) and
\(\sigma_{ar}(B)=\int_{B\cap C_r}p_a(e)\,H(de)\). Using the empirical cell measures in \cref{obj:def:external-projection-handle}, set
\[
D_Y=\sum_{a,r}d_{[0,1]}(\widehat\lambda_{ar},\lambda_{ar}),
\qquad
D_E=\sum_{a,r}d_{\mathcal E}(\widehat\sigma_{ar},\sigma_{ar}),
\qquad
\Lambda_\varepsilon=\varepsilon^{-1}+\varepsilon^{-2}.
\]
For a retained sieve array, the triangle inequality places its outcome and score arrays within \(r_n+D_Y\) and \(r_m+D_E\), respectively, of the population arrays. The quantile-transport arm endpoints are Lipschitz in these distances with coefficient \(\Lambda_\varepsilon\): the reciprocal arm probabilities are bounded by \(\varepsilon^{-1}\), their variation across scores is bounded by \(\varepsilon^{-2}\), and integration of the corresponding distribution-function differences gives the two distances above. Hence each of the four arm endpoints differs from its population endpoint by at most
\[
\delta=\Lambda_\varepsilon(r_n+D_Y+r_m+D_E).
\]
The lower and upper ATE endpoints each combine two arm endpoints. Taking their convex hull and clipping to \([-1,1]\) therefore gives, on the projection branch of \cref{obj:def:external-projection-handle}, for its interval \(C_{n,m}\),
\[
\bigl(\operatorname{len}C_{n,m}
       -\operatorname{len}I(P_{\mathrm{rel}};H,g)\bigr)_+
\le 4\Lambda_\varepsilon(r_n+D_Y+r_m+D_E).
\]
The same inequality holds when the totalized interval in \cref{obj:def:external-projection-handle} is \(\{0\}\), since its length is zero. The empirical-CDF variance bound, integrated over each support interval and summed over the \(2J\) cells, yields
\[
\mathbb E D_Y\le\frac{2J}{\sqrt n},
\qquad
\mathbb E D_E\le\frac{J(2-2\varepsilon)}{\sqrt m}.
\]
The radii \(r_n=4J/(\alpha\sqrt n)\) and
\(r_m=2J(2-2\varepsilon)/(\alpha\sqrt m)\) make the projection interval honest: Markov’s inequality assigns probability at most \(\alpha/2\) to each failed empirical ball, and approximation by compatible rational arrays places the population sharp interval in the resulting closed projection hull on the complementary event. Thus \cref{obj:def:external-excess-risk} and the preceding bounds give
\[
\mathcal R^\star_{n,m,\alpha}(g)
\le 4\Lambda_\varepsilon
\left[
 \left(\frac{4J}{\alpha}+2J\right)n^{-1/2}
 +\left(\frac{2J(2-2\varepsilon)}{\alpha}
          +J(2-2\varepsilon)\right)m^{-1/2}
\right].
\]
In particular, with
\[
M=4\Lambda_\varepsilon
\left(\frac{4J}{\alpha}+2J+
      \frac{2J(2-2\varepsilon)}{\alpha}
      +J(2-2\varepsilon)\right)+1,
\]
we have \(T_{n,m,\alpha}(g)\le M\) for all \(n,m\ge1\), by the definition of \(b_{n,m}\) in \cref{obj:def:external-normalized-risk}.
% lean: CausalSmith.PartialID.UnlinkedPropensityAte.externalExcessRisk_normalized_upper

Fix an ordered triple \(e_1<e_2<e_3\) in one fiber, and define
\[
\mathbf h=(e_3-e_2,\;-(e_3-e_1),\;e_2-e_1),
\quad
u_m=\frac{\sqrt{\log L_\alpha}}
          {\sqrt3\,\|\mathbf h\|_2\sqrt m},
\quad
q=\frac{e_1+e_2+e_3}{3},
\quad
t=\frac12\left\{\frac{e_1}{3}
       +\min\!\left(\frac{e_1+e_2}{3},\frac{e_3}{3}\right)\right\},
\quad s=\frac tq.
\]
The identities \(\sum_i h_i=\sum_i e_i h_i=0\) show that
\[
H_0=\frac13\sum_{i=1}^3\delta_{e_i},
\qquad
H_m=\sum_{i=1}^3\left(\frac13+u_mh_i\right)\delta_{e_i}
\]
are probability laws with the same score mean \(q\) for all sufficiently large \(m\). The strict inequalities
\[
\frac{e_1}{3}<t<
\min\!\left(\frac{e_1+e_2}{3},\frac{e_3}{3}\right),
\qquad 0<s<1,
\]
hold by the choice of \(t\). Writing \(w_{i,m}=1/3+u_mh_i\) for \(i=1,2,3\), continuity and \(u_m\to0\) give, for all sufficiently large \(m\),
\[
w_{i,m}>0\quad(i=1,2,3),\qquad
e_1w_{1,m}<t<
\min\{e_1w_{1,m}+e_2w_{2,m},\,e_3w_{3,m}\}.
\] Under either score law, take \(Y_0=0\), an independent Bernoulli\((s)\) potential outcome \(Y_1\), score-based treatment assignment, and the fixed label of the triple. These constructions satisfy \cref{obj:def:external-law-class}; their released trial laws coincide, with treatment probability \(q\) and treated-success probability \(t\).

For score masses \(w_i\) in this neighborhood, the upper rearrangement of the treated successes fills score \(e_1\) and then score \(e_2\), while the lower rearrangement fits them at score \(e_3\). The endpoint formulas in \cref{obj:def:mean-endpoints} therefore reduce to
\[
U(w)=w_1+\frac{t-e_1w_1}{e_2},
\qquad L(w)=\frac{t}{e_3}.
\]
The control endpoints are zero. The two sharp intervals have common lower endpoint \(L(w)\), and their upper endpoints differ by
\[
\Delta_m
=U\!\left(\tfrac13+u_mh_1,\tfrac13+u_mh_2,\tfrac13+u_mh_3\right)
 -U\!\left(\tfrac13,\tfrac13,\tfrac13\right)
=u_m(e_3-e_2)\left(1-\frac{e_1}{e_2}\right)>0.
\]
By the endpoint-attainment assertion of \cref{obj:lem:fixed-released-law-baseline}, choose causal populations attaining the lower endpoint under \(H_0\) and the upper endpoint under \(H_m\), while preserving their respective released laws. If \(W_0\) denotes the sharp interval length under \(H_0\), their ATEs differ by \(W_0+\Delta_m\).

Only the score logs distinguish these endpoint-attaining experiments. Direct calculation gives
\[
\chi^2(H_m,H_0)
=3u_m^2\|\mathbf h\|_2^2
=\frac{\log L_\alpha}{m}.
\]
For product laws, \(1+\chi^2\) raises to the sample-size power; Cauchy--Schwarz then bounds total variation by half the square root of \(\chi^2\). Using \((1+x/m)^m\le e^x\), the joint experiments consequently satisfy, uniformly in \(n\),
\[
\operatorname{TV}(\mathbb Q_0^{n,m},\mathbb Q_m^{n,m})
\le\frac12
 \sqrt{\left(1+\frac{\log L_\alpha}{m}\right)^m-1}
\le\frac12\sqrt{L_\alpha-1}
=\frac{\beta}{2}.
\]
For any honest interval, transfer coverage of the second ATE to the first experiment and intersect it with coverage of the first ATE. The intersection has first-experiment probability at least
\(1-2\alpha-\operatorname{TV}\ge\beta/2\). On it the interval length is at least \(W_0+\Delta_m\). Taking the supremum over laws and the infimum over honest procedures gives, for every \(n\ge1\) and all sufficiently large \(m\),
\[
\mathcal R^\star_{n,m,\alpha}(g)
\ge\frac{\beta\Delta_m}{2}
=\frac{\beta\sqrt{\log L_\alpha}}{2\sqrt3\,\sqrt m}
  \frac{(e_3-e_2)(1-e_1/e_2)}
       {D(e_1,e_2,e_3)}.
\]
% lean: CausalSmith.PartialID.UnlinkedPropensityAte.scoreLogFixedTriple_eventually_externalExcessRisk_lower

To pass from a fixed triple to \(S_g\), take any number below
\(\beta\sqrt{\log L_\alpha}S_g/(2\sqrt3)\). The definition of the supremum supplies a triple whose displayed coefficient exceeds that number. Its eventual lower bound holds uniformly over \(n\ge1\). The infimum over \(n\) is finite: it is bounded below by this eventual certificate and above by the choice \(n=m\), for which \(b_{m,m}=\sqrt m/2\) and the bound \(T_{m,m,\alpha}(g)\le M\) gives
\(\sqrt m\,\mathcal R^\star_{m,m,\alpha}(g)\le2M\). Taking the liminf and then letting the chosen number increase to the supremum proves
\[
\liminf_{m\to\infty}\inf_{n\ge1}
 \sqrt m\,\mathcal R^\star_{n,m,\alpha}(g)
\ge c_{\mathrm{log}}(g,\alpha).
\]
% lean: CausalSmith.PartialID.UnlinkedPropensityAte.scoreLog_logCertificate_liminf

For the trial certificate, take \(H=\delta_{1/2}\), which is supported on \(\mathcal E\) by \cref{obj:ass:overlap}. Let \(P_0\) have \(Y_0=0\) and an independent Bernoulli\((1/2)\) outcome \(Y_1\), and let \(P_1\) replace its success probability by \(1/2+v_n\), where
\[
v_n=\frac12\sqrt{\frac{\log L_\alpha}{n}}.
\]
Both laws use treatment probability \(E=1/2\), consistency, and the fixed release \(g(1/2)\). Since \(0<\beta<1\) and \(n\ge1\), \(0<v_n<1/2\); both pairs belong to \cref{obj:def:external-law-class}. The fixed score makes the reciprocal treatment weights constant, so \cref{obj:def:mean-endpoints,obj:def:sharp-ate-set} give singleton sharp intervals at \(1/2\) and \(1/2+v_n\).

For a latent Bernoulli outcome, the chi-squared divergence between these two success probabilities is \(4v_n^2\). Adding the common randomized arm, releasing the trial row, and adding the common score log cannot increase total variation. Hence
\[
\operatorname{TV}(\mathbb Q_0^{n,m},\mathbb Q_1^{n,m})
\le\frac12\sqrt{(1+4v_n^2)^n-1}
\le\frac12\sqrt{e^{\log L_\alpha}-1}
=\frac{\beta}{2},
\]
where the middle inequality uses
\((1+\log L_\alpha/n)^n\le L_\alpha\). Under the first experiment, honesty and transfer of the second coverage event put both ATEs in any honest interval with probability at least \(\beta/2\). Its length on that event is at least \(v_n\), while the first sharp interval has length zero. Thus every honest procedure has worst-case expected excess length at least
\[
\frac{\beta v_n}{2}
=\frac{\beta\sqrt{\log L_\alpha}}{4\sqrt n}
=\frac{c_{\mathrm{tr}}(\alpha)}{\sqrt n}.
\]
The full interval \([-1,1]\) is honest, so the infimum defining the minimax risk is over a nonempty set; every candidate satisfies the displayed bound. This proves the finite-sample assertion.
% lean: CausalSmith.PartialID.UnlinkedPropensityAte.trialCenter_honest_expectedLength

Finally, for positive \(n,m\), multiplying the trial bound by \(b_{n,m}\) gives
\[
T_{n,m,\alpha}(g)
\ge c_{\mathrm{tr}}(\alpha)
 \frac{\sqrt m}{\sqrt n+\sqrt m}.
\]
If \(n,m\to\infty\) with \(n/m\to\rho>0\), the fraction tends to
\(1/(1+\sqrt\rho)\). For the log bound, let any
\(\gamma<c_{\mathrm{log}}(g,\alpha)\) be fixed. The preceding liminf inequality implies, eventually in \(m\), uniformly over \(n\ge1\),
\(\mathcal R^\star_{n,m,\alpha}(g)\ge\gamma/\sqrt m\). Hence eventually
\[
T_{n,m,\alpha}(g)
\ge \gamma\frac{\sqrt n}{\sqrt n+\sqrt m}
\longrightarrow
\gamma\frac{\sqrt\rho}{1+\sqrt\rho}.
\]
The uniform bound \(T_{n,m,\alpha}(g)\le M\) makes the normalized liminf finite. Letting \(\gamma\) increase to \(c_{\mathrm{log}}(g,\alpha)\) and combining the two lower bounds yields
\[
\liminf_{\substack{n,m\to\infty\\ n/m\to\rho}}
T_{n,m,\alpha}(g)
\ge
\max\left\{
\frac{c_{\mathrm{tr}}(\alpha)}{1+\sqrt\rho},
\frac{c_{\mathrm{log}}(g,\alpha)\sqrt\rho}{1+\sqrt\rho}
\right\}>0,
\]
as required.
% lean: CausalSmith.PartialID.UnlinkedPropensityAte.externalExcessRisk_lower_certificates
\end{proof}

The trial bound applies at every pair of positive sample sizes. The score-log bound persists when trial size is unrestricted, and the final bound records both contributions at a fixed positive sample-size ratio. These bounds provide the lower-risk component of \cref{obj:thm:external-score-root-order}.

\subsection{Verification scope}

The displayed theorem and lemma declarations and their rendered proofs are machine-checked with Lean \texttt{4.33.0} at source commit \texttt{9135ac11b7d3e6ebd997f8c56bc48396dde396e4}. The machine-readable manifest \texttt{verification\_contract.json} maps each paper object identifier to its Lean declaration, source file, statement and proof audits, external dependencies, and sorry-free status. It records 44 displayed objects and 82 referenced Lean declarations. Objects marked \texttt{matched} have a checked formal counterpart; synthesized presentation definitions organize notation and are identified separately in that manifest. The fixed-law result is credited to \citet[Theorem 3.2 and eq. (3.1)]{FanShermanShum2014} as prior mathematics, while \cref{obj:lem:fixed-released-law-baseline} has no external formal dependency and its present proof is checked under the displayed assumptions.


\section{Proofs of the main results}\label{sec:deferred-proofs}

\begin{proof}[Proof of \cref{obj:prop:compatibility}]
Suppose \(P\in\mathfrak M(H,g,P_{\mathrm{rel}})\). For each label \(r\), score-based assignment, deterministic release, and the score marginal give the treated-arm identity
\[
\begin{aligned}
P_{\mathrm{rel}}(A=1,R=r)
&=P(A=1,g(E)=r)\\
&=\mathbb E_P\!\left[\mathbf 1\{g(E)=r\}p_1(E)\right]
=\int_{C_r}e\,H(de)=q_{1r}.
\end{aligned}
\]
This is the treated-arm instance of \cref{obj:lem:released-arm-cell-masses}; the arm-cell quantities are defined in \cref{obj:def:arm-cell-primitives}. The score marginal and deterministic release also give
\[
P_{\mathrm{rel}}(R=r)=P(g(E)=r)=H(C_r).
\]
Partitioning the label event by arm therefore yields the control-arm identity
\[
\begin{aligned}
P_{\mathrm{rel}}(A=0,R=r)
&=P_{\mathrm{rel}}(R=r)-P_{\mathrm{rel}}(A=1,R=r)\\
&=H(C_r)-\int_{C_r}e\,H(de)
=\int_{C_r}(1-e)\,H(de)=q_{0r}.
\end{aligned}
\]
Thus \(P_{\mathrm{rel}}(A=a,R=r)=q_{ar}\) for both arms.
% lean: CausalSmith.PartialID.UnlinkedPropensityAte.compatibleCellMasses_of_compatibleCausalLaw
Since \(p_0(e)+p_1(e)=1\), summing the arm masses gives
\[
P_{\mathrm{rel}}(R=r)
=\sum_{a\in\mathcal A}P_{\mathrm{rel}}(A=a,R=r)
=\sum_{a\in\mathcal A}q_{ar}
=\int_{C_r}1\,H(de)
=H(C_r).
\]
The treated-arm identity above is the second required equality.
% lean: CausalSmith.PartialID.UnlinkedPropensityAte.releasedLabelMass_sum_arms
% lean: CausalSmith.PartialID.UnlinkedPropensityAte.compatibility_armCellMass_sum

Conversely, suppose the two stated equalities hold for every \(r\). Partitioning each label event by arm and using \(p_0=1-p_1\) yields
\[
\begin{aligned}
P_{\mathrm{rel}}(A=0,R=r)
&=P_{\mathrm{rel}}(R=r)-P_{\mathrm{rel}}(A=1,R=r)\\
&=H(C_r)-\int_{C_r}e\,H(de)
=\int_{C_r}(1-e)\,H(de)=q_{0r}.
\end{aligned}
\]
The assumed treated equality is \(P_{\mathrm{rel}}(A=1,R=r)=q_{1r}\). Thus \(P_{\mathrm{rel}}(A=a,R=r)=q_{ar}\) for both arms and every label.
% lean: CausalSmith.PartialID.UnlinkedPropensityAte.compatibleCellMasses_of_label_treated

For every \((a,r)\in\mathcal A\times\mathcal R_J\), define a probability law on the outcome interval by
\[
F^\star_{ar}(\mathcal U):=
\begin{cases}
P_{\mathrm{rel}}(Y\in\mathcal U\mid A=a,R=r), & P_{\mathrm{rel}}(A=a,R=r)>0,\\
\mathbf 1\{0\in\mathcal U\}, & P_{\mathrm{rel}}(A=a,R=r)=0,
\end{cases}
\qquad \mathcal U\subseteq[0,1]\text{ Borel}.
\]
On positive-mass pairs this is \(F_{ar}\) from \cref{obj:def:arm-cell-primitives}; on zero-mass pairs it is point mass at zero. In particular, for every Borel outcome set \(\mathcal U\),
\[
P_{\mathrm{rel}}(Y\in\mathcal U,A=a,R=r)
=P_{\mathrm{rel}}(A=a,R=r)F^\star_{ar}(\mathcal U)
=q_{ar}F^\star_{ar}(\mathcal U).
\]
Define a law for the latent variables by
\[
\Lambda(de,dy_0,dy_1)
=H(de)\,F^\star_{0,g(e)}(dy_0)\,F^\star_{1,g(e)}(dy_1).
\]
Its score marginal is \(H\), and its two potential outcomes are conditionally independent given the score. The finite label alphabet and measurability of \(g\) make these kernels measurable.
% lean: CausalSmith.PartialID.UnlinkedPropensityAte.reconstructedLatentLaw

From \(\Lambda\), construct \(\widetilde P\) by drawing \(A\) conditionally on \((E,Y_0,Y_1)\) with probabilities \(p_a(E)\), then setting \(R=g(E)\) and \(Y=AY_1+(1-A)Y_0\). Equivalently, in the full-row coordinate order \((E,Y_0,Y_1,A,Y,R)\), its law is
\[
\widetilde P
=\sum_{a\in\mathcal A}
\bigl((e,y_0,y_1)\mapsto(e,y_0,y_1,a,ay_1+(1-a)y_0,g(e))\bigr)_{\#}
\bigl(p_a(e)\,\Lambda(de,dy_0,dy_1)\bigr).
\] Overlap gives \(0\leq p_a(e)\leq1\), and
\[
p_0(e)+p_1(e)=1,\qquad
\widetilde P(E\in\mathcal S)
=\int_{\mathcal S}\sum_{a\in\mathcal A}p_a(e)\,H(de)
=H(\mathcal S)
\]
for every Borel score set \(\mathcal S\). Hence \(\widetilde P\) is a probability law with score marginal \(H\). Its conditional assignment probabilities are \(p_a(E)\), while consistency and deterministic release hold by construction.
% lean: CausalSmith.PartialID.UnlinkedPropensityAte.reconstructedFullLaw
% lean: CausalSmith.PartialID.UnlinkedPropensityAte.reconstructedFullLaw_isProbabilityMeasure
% lean: CausalSmith.PartialID.UnlinkedPropensityAte.reconstructedFullLaw_scoreMarginal
% lean: CausalSmith.PartialID.UnlinkedPropensityAte.reconstructedFullLaw_randomizedAssignment
% lean: CausalSmith.PartialID.UnlinkedPropensityAte.reconstructedFullLaw_consistency
% lean: CausalSmith.PartialID.UnlinkedPropensityAte.reconstructedFullLaw_deterministicRelease

Finally, for every \(a,r,\mathcal U\) as above,
\[
\begin{aligned}
\widetilde P(R=r,A=a,Y\in\mathcal U)
&=\int_{C_r}p_a(e)F^\star_{ar}(\mathcal U)\,H(de)\\
&=q_{ar}F^\star_{ar}(\mathcal U)
=P_{\mathrm{rel}}(R=r,A=a,Y\in\mathcal U).
\end{aligned}
\]
These events determine the law of \((R,A,Y)\), including zero-mass cells. Thus its law under \(\widetilde P\) is \(P_{\mathrm{rel}}\), and \(\widetilde P\in\mathfrak M(H,g,P_{\mathrm{rel}})\).
% lean: CausalSmith.PartialID.UnlinkedPropensityAte.reconstructedFullLaw_releasedLaw
% lean: CausalSmith.PartialID.UnlinkedPropensityAte.reconstructedFullLaw_compatibleCausalLaw
\end{proof}

\begin{proof}[Proof of \cref{obj:thm:all-label-ambiguity}]
\textit{Step 1 (general upper bound).}
For each arm and label, define
\[
\Gamma_{ar}
:=\inf_{c\in[1/(1-\varepsilon),\,1/\varepsilon]}
       \int_{C_r}|1-cp_a(e)|\,H(de),
\qquad
\Gamma_{\mathrm{all}}:=\sum_{r\in\mathcal R_J}
       \bigl(\Gamma_{1r}+\Gamma_{0r}\bigr).
\]
Consider any causal law \(\widetilde P\) in the supremum defining \(D_H(g)\), and set
\[
\widetilde P_{\mathrm{rel}}
:=\mathcal L_{\widetilde P}(R,A,Y).
\]
For a cell with \(q_{ar}>0\), let \(\widetilde F_{ar}\) be its observed outcome law under \(\widetilde P_{\mathrm{rel}}\), and define its contribution to the interval width by
\[
\Delta_{ar}(\widetilde P_{\mathrm{rel}})
:=q_{ar}\int_0^1 Q_{\widetilde F_{ar}}(u)
       \bigl\{Q_{W_{ar}}(u)-Q_{W_{ar}}(1-u)\bigr\}\,du.
\]
Set \(\Delta_{ar}(\widetilde P_{\mathrm{rel}})=0\) when \(q_{ar}=0\). By \cref{obj:lem:fixed-released-law-baseline} and the definitions of the mean endpoints,
\[
\operatorname{len}I(\widetilde P_{\mathrm{rel}};H,g)
=\sum_{r\in\mathcal R_J}
 \bigl\{\Delta_{1r}(\widetilde P_{\mathrm{rel}})
       +\Delta_{0r}(\widetilde P_{\mathrm{rel}})\bigr\}.
\]

Fix a positive-mass cell and \(c\in[1/(1-\varepsilon),1/\varepsilon]\). Since \(0\le Q_{\widetilde F_{ar}}(u)\le1\) almost everywhere, the positive and negative parts give
\[
\begin{aligned}
\int_0^1 Q_{\widetilde F_{ar}}(u)Q_{W_{ar}}(u)\,du
&\le c\int_0^1Q_{\widetilde F_{ar}}(u)\,du
   +\int_0^1\bigl(Q_{W_{ar}}(u)-c\bigr)_+\,du,\\
\int_0^1 Q_{\widetilde F_{ar}}(u)Q_{W_{ar}}(1-u)\,du
&\ge c\int_0^1Q_{\widetilde F_{ar}}(u)\,du
   -\int_0^1\bigl(c-Q_{W_{ar}}(1-u)\bigr)_+\,du.
\end{aligned}
\]
The quantiles at \(u\) and \(1-u\) each have law \(W_{ar}\) under uniform \(u\). Subtraction therefore yields
\[
\frac{\Delta_{ar}(\widetilde P_{\mathrm{rel}})}{q_{ar}}
\le \mathbb E|W_{ar}-c|.
\]
The defining score law of \(E_{ar}\) and \(p_a(e)>0\) give the scaling identity
\[
q_{ar}\mathbb E|W_{ar}-c|
=\int_{C_r}p_a(e)\left|\frac1{p_a(e)}-c\right|H(de)
=\int_{C_r}|1-cp_a(e)|\,H(de).
\]
Taking the infimum shows \(\Delta_{ar}(\widetilde P_{\mathrm{rel}})\le\Gamma_{ar}\). If \(q_{ar}=0\), then \(H(C_r)=0\), since \(p_a(e)\ge\varepsilon>0\) on \(\mathcal E\); hence both sides are zero. Summing over all cells proves
\[
\operatorname{len}I(\widetilde P_{\mathrm{rel}};H,g)
\le\Gamma_{\mathrm{all}},
\qquad D_H(g)\le\Gamma_{\mathrm{all}}.
\]
% lean: CausalSmith.PartialID.UnlinkedPropensityAte.worstCaseAmbiguity_le_sum_cellAbsoluteDeviation

\textit{Step 2 (width of the half-outcome release).}
For the \(P_{\mathrm{rel}}\) in the statement, every positive-mass cell has outcome law Bernoulli\((1/2)\), whose quantile is zero below \(1/2\) and one above \(1/2\), up to a null set. Put \(m_{ar}:=Q_{W_{ar}}(1/2)\) in such a cell. Monotonicity of the quantile and a change of variable give
\[
\begin{aligned}
\frac{\Delta_{ar}(P_{\mathrm{rel}})}{q_{ar}}
&=\int_{1/2}^1Q_{W_{ar}}(u)\,du
  -\int_0^{1/2}Q_{W_{ar}}(u)\,du\\
&=\int_0^1|Q_{W_{ar}}(u)-m_{ar}|\,du
 =\mathbb E|W_{ar}-m_{ar}|.
\end{aligned}
\]
Overlap places \(W_{ar}\), and thus its median \(m_{ar}\), in
\([1/(1-\varepsilon),1/\varepsilon]\). The preceding cellwise bound applies to this fair outcome law for every admissible \(c\), while \(m_{ar}\) attains that bound. The scaling identity consequently gives
\[
\Delta_{ar}(P_{\mathrm{rel}})
=q_{ar}\inf_{c\in[1/(1-\varepsilon),\,1/\varepsilon]}
       \mathbb E|W_{ar}-c|
=\Gamma_{ar}.
\]
The equality also holds in zero-mass cells by the zero-mass argument above. Summing the armwise contributions yields
\[
\operatorname{len}I(P_{\mathrm{rel}};H,g)
=\Gamma_{\mathrm{all}}.
\]
% lean: CausalSmith.PartialID.UnlinkedPropensityAte.halfOutcome_sharpATELength_eq_sum_cellAbsoluteDeviation

\textit{Step 3 (compatible laws and endpoints).}
The cell masses satisfy
\[
\sum_{r\in\mathcal R_J}\sum_{a\in\{0,1\}}q_{ar}
=\int_{\mathcal E}\bigl(p_1(e)+p_0(e)\bigr)\,H(de)=1.
\]
Thus \(P_{\mathrm{rel}}\) is a probability law. For every label \(r\), its defining mixture gives
\[
P_{\mathrm{rel}}(R=r)=q_{1r}+q_{0r}=H(C_r),
\qquad
P_{\mathrm{rel}}(A=1,R=r)=q_{1r}
=\int_{C_r}e\,H(de).
\]
By \cref{obj:prop:compatibility}, there is a compatible full law producing \(P_{\mathrm{rel}}\). Applying \cref{obj:lem:fixed-released-law-baseline} to this released law gives compatible probability laws \(P_-,P_+\) with
\[
\mathbb E_{P_-}[Y_1-Y_0]=\underline\mu_1-\overline\mu_0,
\qquad
\mathbb E_{P_+}[Y_1-Y_0]=\overline\mu_1-\underline\mu_0,
\]
and hence
\[
\mathbb E_{P_+}[Y_1-Y_0]-\mathbb E_{P_-}[Y_1-Y_0]
=\operatorname{len}I(P_{\mathrm{rel}};H,g)
=\Gamma_{\mathrm{all}}.
\]
% lean: CausalSmith.PartialID.UnlinkedPropensityAte.halfOutcome_endpoint_attainment

\textit{Step 4 (supremum and conclusion).}
The set of widths in \cref{obj:def:worst-case-ambiguity} is bounded above by \(\Gamma_{\mathrm{all}}\), by the general upper bound. The compatible law constructed for \(P_{\mathrm{rel}}\) puts its width in that set, and the half-outcome calculation shows that this width equals \(\Gamma_{\mathrm{all}}\). Therefore
\[
D_H(g)=\Gamma_{\mathrm{all}}
=\sum_{r\in\mathcal R_J}\left[
\inf_{c\in[1/(1-\varepsilon),\,1/\varepsilon]}
 \int_{C_r}|1-ce|\,H(de)
+
\inf_{c\in[1/(1-\varepsilon),\,1/\varepsilon]}
 \int_{C_r}|1-c(1-e)|\,H(de)
\right].
\]
Substituting this equality into the attained endpoint gap above proves the asserted gap. The hypotheses also cover the index-set boundary: \(\mathcal E\) contains \(\varepsilon\), so the value \(g(\varepsilon)\) ensures \(\mathcal R_J\) is nonempty.
% lean: CausalSmith.PartialID.UnlinkedPropensityAte.worstCaseAmbiguity_eq
\end{proof}

\begin{proof}[Proof of \cref{obj:thm:universal-point-identification}]
\begin{enumerate}
\item For \(a\in\{0,1\}\) and \(r\in\mathcal R_J\), set
\[
\delta_{ar}:=
\inf_{c\in[1/(1-\varepsilon),\,1/\varepsilon]}
\int_{C_r}|1-cp_a(e)|\,H(de),
\qquad p_1(e)=e,\quad p_0(e)=1-e.
\]
The coefficient interval is nonempty under overlap. Each defining integral is finite and nonnegative, so \(\delta_{ar}\geq0\). By \cref{obj:thm:all-label-ambiguity},
\[
D_H(g)=\sum_{r\in\mathcal R_J}(\delta_{1r}+\delta_{0r}).
\]
Thus \(D_H(g)=0\) implies \(\delta_{1r}=0\) for every \(r\): each term in this finite sum is nonnegative. % lean: CausalSmith.PartialID.UnlinkedPropensityAte.worstCaseAmbiguity_eq

\item Fix \(r\) with \(H(C_r)>0\), and define, for every real \(t\),
\[
\Phi_r(t):=\int_{C_r}|1-te|\,H(de).
\]
Since \(0\leq e\leq1\) on \(\mathcal E\), for all real \(t,t'\),
\[
|\Phi_r(t)-\Phi_r(t')|
\leq \int_{C_r}|t-t'|e\,H(de)
\leq |t-t'|H(C_r)
\leq |t-t'|.
\]
Hence \(\Phi_r\) is continuous and attains its minimum on the compact coefficient interval. Let \(c_r\) attain it. The preceding step gives
\[
0=\delta_{1r}=\Phi_r(c_r)
=\int_{C_r}|1-c_re|\,H(de),
\]
so \(1-c_re=0\) for \(H\)-almost every \(e\in C_r\). Because \(H(C_r)>0\), there is an \(x_r\in C_r\) with \(1-c_rx_r=0\). Also \(c_r\geq1/(1-\varepsilon)>0\). Consequently,
\[
c_r(e-x_r)=0
\quad\text{for \(H\)-almost every \(e\in C_r\),}
\]
and therefore \(e=x_r\) there almost surely. This proves the cellwise condition. % lean: CausalSmith.PartialID.UnlinkedPropensityAte.cellAbsoluteDeviation_minimizer

\item Conversely, suppose the cellwise condition holds. Choose its value \(x_r\) on each positive-mass cell and define the map on the entire label set by
\[
h(r):=
\begin{cases}
x_r,&H(C_r)>0,\\
\varepsilon,&H(C_r)=0,
\end{cases}
\qquad r\in\mathcal R_J.
\]
Here \(\varepsilon\in\mathcal E\) by overlap. The finite label set makes \(h\) measurable. On every positive-mass cell, \(e=h(r)\) almost surely; the same assertion is vacuous on a zero-mass cell. Taking the finite union of these cellwise assertions yields
\[
e=h(g(e))\qquad\text{for \(H\)-almost every \(e\)}.
\]
% lean: CausalSmith.PartialID.UnlinkedPropensityAte.scoreRecovery_iff_cellwiseAE

\item For every \(a\in\{0,1\}\) and \(r\in\mathcal R_J\), set
\[
c_{ar}:=\frac{1}{p_a(h(r))}.
\]
Overlap gives
\[
\varepsilon\leq p_a(h(r))\leq1-\varepsilon,
\qquad
\frac1{1-\varepsilon}\leq c_{ar}\leq\frac1\varepsilon.
\]
The cellwise equality established above gives \(p_a(e)=p_a(h(r))\) almost surely on \(C_r\), including when that cell has zero mass. Thus
\[
0\leq\delta_{ar}
\leq\int_{C_r}|1-c_{ar}p_a(e)|\,H(de)=0.
\]
Every summand in the all-label formula is zero, and hence \(D_H(g)=0\). Finally, if a measurable \(h\) satisfies \(e=h(g(e))\) \(H\)-almost surely, then on each positive-mass \(C_r\) its value \(h(r)\) witnesses the cellwise condition. The two stated conditions are therefore each equivalent to \(D_H(g)=0\). % lean: CausalSmith.PartialID.UnlinkedPropensityAte.worstCaseAmbiguity_eq_zero_iff
\end{enumerate}
\end{proof}

\begin{proof}[Proof of \cref{obj:thm:k-label-rate}]
The overlap condition gives \(\ell>0\). First fix \(H\in\mathcal H^{\mathrm{lb}}_{\varepsilon,m_f}\) and an arbitrary measurable \(g\in\mathcal G_K\). For each label \(r\), set
\[
C_r=\{e\in\mathcal E:g(e)=r\},
\qquad h_r=\operatorname{Leb}(C_r).
\]
The measurable cells partition \(\mathcal E\), including any empty cells, so \(h_r\geq0\) and \(\sum_r h_r=\ell\). % lean: realScoreCell_intervalPartition

For every \(\zeta\in\mathbb R\) and \(t\geq0\), the points within distance \(t\) of \(\zeta\) occupy an interval of length \(2t\). Thus
\[
\operatorname{Leb}\{e\in C_r:|e-\zeta|>t\}\geq h_r-2t,
\]
and the layer-cake identity gives
\[
\int_{C_r}|e-\zeta|\,de
\geq\int_0^{h_r/2}(h_r-2t)\,dt
=\frac{h_r^2}{4}.
\] % lean: boundedCell_absDeviation_integral_lower

Since \(H(de)=f(e)\,de\) with \(f\geq m_f\) almost everywhere, it follows for every \(\zeta\in\mathbb R\) that
\[
\int_{C_r}|e-\zeta|\,H(de)
\geq m_f\int_{C_r}|e-\zeta|\,de
\geq \frac{m_fh_r^2}{4}.
\] % lean: scoreCell_absDeviation_lower

Write \(p_1(e)=e\) and \(p_0(e)=1-e\), as in \cref{obj:thm:all-label-ambiguity}. For any \(c\in[1/(1-\varepsilon),1/\varepsilon]\), define the two score centers by
\[
z_{1,c}=c^{-1},\qquad z_{0,c}=1-c^{-1}.
\]
Here \(c>0\), and direct algebra gives
\[
|1-cp_a(e)|=c|e-z_{a,c}|,
\qquad a\in\{0,1\}.
\] % lean: reciprocalDeviation_eq_scoreDistance

Because \(c\geq1/(1-\varepsilon)\), the preceding cell bound yields, for either arm,
\[
\int_{C_r}|1-cp_a(e)|\,H(de)
\geq \frac{m_fh_r^2}{4(1-\varepsilon)}.
\] % lean: cellAbsoluteDeviation_lower_cellLengthSq

By the exact cell formula in \cref{obj:thm:all-label-ambiguity},
\[
D_H(g)
=\sum_r\sum_{a\in\{0,1\}}
\inf_{c\in[1/(1-\varepsilon),\,1/\varepsilon]}
\int_{C_r}|1-cp_a(e)|\,H(de)
\geq \frac{m_f}{2(1-\varepsilon)}\sum_r h_r^2.
\] % lean: worstCaseAmbiguity_eq

There are \(K\) labels, so Cauchy–Schwarz gives
\[
\sum_r h_r^2
\geq \frac{(\sum_r h_r)^2}{K}
=\frac{\ell^2}{K}.
\] % lean: sum_cellMass_sq_lower
Consequently every \(g\in\mathcal G_K\) satisfies
\[
D_H(g)\geq\frac{m_f\ell^2}{2(1-\varepsilon)K}.
\] % lean: worstCaseAmbiguity_finiteLabel_lower
The class \(\mathcal G_K\) is nonempty because \(K\geq1\) permits a constant measurable release. Taking the infimum as in \cref{obj:def:optimal-k-ambiguity} proves the lower bound for \(\Delta_K(H)\). % lean: optimalKAmbiguity_finiteLabel_lower

Now let \(H\) be any probability law on \(\mathcal E\), and take the equal-width release \(g\) in the statement. It is measurable: the score transformation inside the floor is continuous, the floor is measurable, and clipping its value to the finite label set preserves measurability. % lean: equalWidthRelease_measurable

Index its labels by \(r=1,\ldots,K\), write \(C_r=\{e\in\mathcal E:g(e)=r\}\), and define
\[
z_r=\varepsilon+\left(r-\frac12\right)\frac{\ell}{K}.
\]
For \(e\in C_r\), the floor rule places \(e\) between \(\varepsilon+(r-1)\ell/K\) and \(\varepsilon+r\ell/K\). At the right endpoint of \(\mathcal E\), clipping assigns the last label, for which the same closed bounds hold. Hence \(z_r\in\mathcal E\) and
\[
|e-z_r|\leq\frac{\ell}{2K}\qquad(e\in C_r).
\] % lean: equalWidthRelease_cell_subset_closedBin

The candidates \(z_r^{-1}\) and \((1-z_r)^{-1}\) lie in \([1/(1-\varepsilon),1/\varepsilon]\). Moreover,
\[
z_r(1-z_r)-\varepsilon(1-\varepsilon)
=(z_r-\varepsilon)(1-\varepsilon-z_r)\geq0.
\]
Thus, for \(e\in C_r\),
\[
\left|1-\frac e{z_r}\right|
+\left|1-\frac{1-e}{1-z_r}\right|
=\frac{|e-z_r|}{z_r(1-z_r)}
\leq \frac{|e-z_r|}{\varepsilon(1-\varepsilon)}.
\] % lean: sum_arm_reciprocalDeviation_le

Using these candidates in the two cell infima, then applying the midpoint bound, gives
\[
\sum_{a\in\{0,1\}}
\inf_{c\in[1/(1-\varepsilon),\,1/\varepsilon]}
\int_{C_r}|1-cp_a(e)|\,H(de)
\leq
\frac{1}{\varepsilon(1-\varepsilon)}
\int_{C_r}|e-z_r|\,H(de)
\leq
\frac{\ell}{2\varepsilon(1-\varepsilon)K}H(C_r).
\] % lean: sum_cellAbsoluteDeviation_le_centeredScoreIntegral

The cells partition \(\mathcal E\) and \(H(\mathcal E)=1\). Summing the cell bounds therefore yields
\[
D_H(g)\leq
\frac{\ell}{2\varepsilon(1-\varepsilon)K}
\sum_r H(C_r)
=
\frac{\ell}{2\varepsilon(1-\varepsilon)K}.
\] % lean: equalWidthRelease_worstCaseAmbiguity_le_all

Finally, the exact cell formula shows \(D_H(g')\geq0\) for every measurable release \(g'\), while the measurable \(g\) above belongs to \(\mathcal G_K\). Hence the defining infimum satisfies
\[
\Delta_K(H)\leq D_H(g)
\leq\frac{\ell}{2\varepsilon(1-\varepsilon)K}.
\] % lean: optimalKAmbiguity_bounds
\end{proof}

\begin{proof}[Proof of \cref{obj:thm:sharp-high-resolution-constant}]
By \cref{obj:ass:overlap}, \(\mathcal E\) is a nondegenerate compact interval and \(e,1-e\geq\varepsilon>0\) throughout \(\mathcal E\). Continuity and strict positivity give
\[
0<\min_{e\in\mathcal E}f(e)
\leq\max_{e\in\mathcal E}f(e)<\infty.
\]
% lean: CausalSmith.PartialID.UnlinkedPropensityAte.exists_boundedScoreDensity_of_continuous

For \(x,z\in\mathcal E\), define the two distortion coefficients
\[
\beta_1(x,z)=\frac{f(x)}{z},
\qquad
\beta_0(x,z)=\frac{f(x)}{1-z}.
\]
Both are continuous and strictly positive on \(\mathcal E^2\). Their diagonal weight is
\[
w(x):=\beta_0(x,x)+\beta_1(x,x)
=\frac{f(x)}{x(1-x)}.
\]
Thus \cref{obj:lem:paired-high-resolution-quantization} applies with two components. In particular, if \(V_K\) denotes its optimal paired cost for these coefficients, then
\[
K V_K\longrightarrow
\frac14\left(\int_{\varepsilon}^{1-\varepsilon}
\sqrt{\frac{f(x)}{x(1-x)}}\,dx\right)^2.
\]
% lean: CausalSmith.PartialID.UnlinkedPropensityAte.pairedQuantization_tendsto
% lean: CausalSmith.PartialID.UnlinkedPropensityAte.highResolutionBeta_diagonal

We identify this paired cost with the release objective. Fix \(K>0\) and a measurable release \(g:\mathcal E\to\{0,\ldots,K-1\}\), and put \(C_r=\{x\in\mathcal E:g(x)=r\}\). For centers \(\mathbf z=(z_{r,a})_{r,a}\in\mathcal E^{\{0,\ldots,K-1\}\times\{0,1\}}\), define
\[
\Phi_K(g,\mathbf z)
:=
\sum_{r=0}^{K-1}\sum_{a=0}^{1}
\int_{C_r}\beta_a(x,z_{r,a})|x-z_{r,a}|\,dx.
\]
The changes of variables \(z_{r,1}=c^{-1}\) and \(z_{r,0}=1-c^{-1}\) each map the coefficient range in \cref{obj:thm:all-label-ambiguity} bijectively onto \(\mathcal E\), and give
\[
f(x)|1-cx|
=\beta_1(x,z_{r,1})|x-z_{r,1}|,
\qquad
f(x)|1-c(1-x)|
=\beta_0(x,z_{r,0})|x-z_{r,0}|.
\]
The cellwise infima are independent. Each has a nonempty set of feasible values bounded below by zero, including when its cell is empty. Hence \cref{obj:thm:all-label-ambiguity} yields
\[
D_H(g)=\inf_{\mathbf z}\Phi_K(g,\mathbf z).
\]
Taking the infimum over releases flattens these independent choices. Their cells are measurable, pairwise disjoint, and cover \(\mathcal E\); conversely, every such \(K\)-cell partition defines a measurable release by assigning its cell indices as labels. Empty cells are allowed in both descriptions. Therefore \cref{obj:def:optimal-k-ambiguity} gives the exact identity
\[
\Delta_K(H)
=\inf_{g\in\mathcal G_K}\inf_{\mathbf z}\Phi_K(g,\mathbf z)
=V_K.
\]
It follows that
\[
K\Delta_K(H)\longrightarrow
\frac14\left(\int_{\varepsilon}^{1-\varepsilon}
\sqrt{\frac{f(x)}{x(1-x)}}\,dx\right)^2.
\]
% lean: CausalSmith.PartialID.UnlinkedPropensityAte.optimalKAmbiguity_eq_optimalPairedCost

To construct the releases, set
\[
A_H:=\int_{\varepsilon}^{1-\varepsilon}\sqrt{w(x)}\,dx.
\]
For each \(K>0\), choose boundaries \(t_{0,K},\ldots,t_{K,K}\) by
\[
t_{0,K}=\varepsilon,\qquad t_{K,K}=1-\varepsilon,\qquad
\int_{\varepsilon}^{t_{j,K}}\sqrt{w(x)}\,dx=\frac{j}{K}A_H
\quad(0\leq j\leq K).
\]
Positivity of \(w\) makes these boundaries strictly increasing. Let \(g_K\) assign label \(j\) to \([t_{j,K},t_{j+1,K})\) for \(j<K-1\), and label \(K-1\) to \([t_{K-1,K},t_{K,K}]\). Put
\[
\zeta_{j,K}:=\frac{t_{j,K}+t_{j+1,K}}2,
\qquad
U_K:=\Phi_K(g_K,\mathbf z)
\quad\text{with }z_{j,0}=z_{j,1}=\zeta_{j,K}.
\]
The construction clause of \cref{obj:lem:paired-high-resolution-quantization} ensures that these cells form a measurable partition, that every midpoint lies in \(\mathcal E\), and that
\[
K U_K\longrightarrow
\frac14 A_H^2.
\]
% lean: CausalSmith.PartialID.UnlinkedPropensityAte.highResolutionOrderedRelease
% lean: CausalSmith.PartialID.UnlinkedPropensityAte.pairedQuantization_tendsto

For every \(K>0\), optimality over releases gives the lower bound, while the chosen centers give a feasible value for each cellwise infimum:
\[
\Delta_K(H)\leq D_H(g_K)
=\inf_{\mathbf z}\Phi_K(g_K,\mathbf z)
\leq U_K.
\]
Multiplication by \(K\geq0\) preserves both inequalities. At \(K=0\), assign zero to the two auxiliary scaled quantities; both bounds then read \(0\leq0\leq0\). The two established limits squeeze \(K D_H(g_K)\) to \(\frac14A_H^2\).
% lean: CausalSmith.PartialID.UnlinkedPropensityAte.optimalKAmbiguity_le_highResolutionOrderedRelease
% lean: CausalSmith.PartialID.UnlinkedPropensityAte.highResolutionOrderedRelease_worstCaseAmbiguity_le_compandingCost
% lean: CausalSmith.PartialID.UnlinkedPropensityAte.tendsto_K_mul_optimalKAmbiguity

Finally, each displayed cell of \(g_K\) is an interval and is measurable. If \(x,z\) carry the same label and \(x\leq y\leq z\), that interval also contains \(y\); thus the label cells have the asserted interval property.
% lean: CausalSmith.PartialID.UnlinkedPropensityAte.highResolutionOrderedRelease_measurable
% lean: CausalSmith.PartialID.UnlinkedPropensityAte.highResolutionOrderedRelease_cell_ordConnected
\end{proof}

\begin{proof}[Proof of \cref{obj:thm:minimax-honest-length}]
Put
\[
\ell=1-2\varepsilon,\qquad \beta=1-2\alpha,\qquad
B_\alpha=\frac{\beta\sqrt{\log(1+\beta^2)}}4,\qquad
U_{\varepsilon,\alpha}
=\frac{\ell}{\varepsilon(1-\varepsilon)}
+\frac8\varepsilon\sqrt{\log(4/\alpha)},\qquad
C_{\varepsilon,\alpha}=\max\{B_\alpha,U_{\varepsilon,\alpha}\}.
\]
The overlap and miscoverage conditions give \(\ell>0\), \(\beta>0\), and \(B_\alpha>0\); hence \(C_{\varepsilon,\alpha}>0\). % lean: CausalSmith.PartialID.UnlinkedPropensityAte.minimaxHonestLength_order

First fix any probability score law \(H\), any \(n,K\geq1\), and any honest pair \((g,C_n)\). Set
\[
u_n=\frac12\sqrt{\frac{\log(1+\beta^2)}{n}}.
\]
Construct two causal laws \(P^{(0)},P^{(1)}\) as follows. Under both, \(E\sim H\), \(Y_0=0\), \(R=g(E)\), and, conditional on the latent variables, \(A\) is Bernoulli with parameter \(E\). Independently of \(E\), take \(Y_1\) to be Bernoulli with parameter \(1/2\) under \(P^{(0)}\) and \(1/2+u_n\) under \(P^{(1)}\); set \(Y=AY_1+(1-A)Y_0\). Since \(0<u_n<1/2\), both laws belong to \(\mathfrak P(H,g)\) of \cref{obj:def:causal-law-class}, and their average treatment effects are \(1/2\) and \(1/2+u_n\). % lean: CausalSmith.PartialID.UnlinkedPropensityAte.honestProcedure_sampling_baseline_expectedLength_lower

Let \(Q_j\) be the released \(n\)-row law under \(P^{(j)}\), for \(j=0,1\). The released rows are measurable functions of Bernoulli potential outcomes and a common, independent score-and-assignment design. Total variation contracts under this readout. The chi-squared divergence of Bernoulli\((1/2+u_n)\) from Bernoulli\((1/2)\) is \(4u_n^2\); its product-law identity and Cauchy–Schwarz therefore give
\[
\operatorname{TV}(Q_0,Q_1)
\leq \frac12\sqrt{(1+4u_n^2)^n-1}.
\]
Writing \(x=\log(1+\beta^2)\), we have \(4u_n^2=x/n\), so \((1+x/n)^n\leq e^x=1+\beta^2\). Consequently,
\[
\operatorname{TV}(Q_0,Q_1)\leq\frac{\beta}{2}.
\]
The two coverage events have probabilities at least \(1-\alpha\) under their respective laws. Transferring the second event to \(Q_0\) and intersecting the events shows that, with \(Q_0\)-probability at least
\(1-2\alpha-\operatorname{TV}(Q_0,Q_1)\geq\beta/2\), the interval contains both treatment effects. Its length is then at least \(u_n\). Thus
\[
\sup_{P\in\mathfrak P(H,g)}
\mathbb E_{P_{\mathrm{rel}}^{\otimes n}}|C_n|
\geq \mathbb E_{Q_0}|C_n|
\geq \frac{\beta u_n}{2}
=B_\alpha n^{-1/2}.
\]
% lean: CausalSmith.PartialID.UnlinkedPropensityAte.honestProcedure_sampling_worstCaseExpectedLength_lower

This bound holds for every honest pair. The collection of such pairs is nonempty: a constant release paired with the interval \([-1,1]\) is honest. Their expected lengths lie in \([0,2]\). Taking the infimum as in \cref{obj:def:minimax-honest-length} yields
\[
B_\alpha n^{-1/2}\leq\mathcal L^\star_{n,K,\alpha}(H).
\]
% lean: CausalSmith.PartialID.UnlinkedPropensityAte.minimaxHonestLength_sampling_lower

Now suppose \(H\in\mathcal H^{\mathrm{lb}}_{\varepsilon,m_f}\), where \(m_f>0\). By \cref{obj:def:score-lower-density-class}, \(H(de)=f(e)\,de\) with \(f\geq m_f\) almost everywhere. Fix an arbitrary \(g\in\mathcal G_K\), and write
\[
C_r=\{e\in\mathcal E:g(e)=r\},\qquad
h_r=\operatorname{Leb}(C_r),\qquad r\in\mathcal R_K.
\]
These measurable cells partition \(\mathcal E\), including when some are empty or disconnected, so \(\sum_rh_r=\ell\). For \(c\in[1/(1-\varepsilon),1/\varepsilon]\), put
\[
z_{1,c}=c^{-1},\qquad z_{0,c}=1-c^{-1}.
\]
Both centers lie in \(\mathcal E\). With \(p_1(e)=e\) and \(p_0(e)=1-e\), as in \cref{obj:thm:all-label-ambiguity}, for either arm \(a\) we have
\[
\int_{C_r}|1-cp_a(e)|\,H(de)
=c\int_{C_r}|e-z_{a,c}|\,H(de)
\geq\frac{m_f}{1-\varepsilon}
       \int_{C_r}|e-z_{a,c}|\,de.
\]
% lean: CausalSmith.PartialID.UnlinkedPropensityAte.cellAbsoluteDeviation_lower_cellLengthSq

Here is the geometric bound used for each cell. For any measurable set \(V\) of finite Lebesgue measure \(h\) and any real \(z\), the interval \(\{|e-z|\leq t\}\) has measure at most \(2t\). The layer-cake identity gives, also when \(h=0\),
\[
\int_V|e-z|\,de
=\int_0^\infty\operatorname{Leb}\{e\in V:|e-z|>t\}\,dt
\geq\int_0^{h/2}(h-2t)\,dt
=\frac{h^2}{4}.
\]
Applying this with \(V=C_r\) bounds each arm’s infimum in the exact ambiguity formula of \cref{obj:thm:all-label-ambiguity}. Since \(\sum_rh_r^2\geq(\sum_rh_r)^2/K=\ell^2/K\), that formula yields
\[
D_H(g)\geq
\frac{m_f}{2(1-\varepsilon)}\sum_r h_r^2
\geq\frac{m_f\ell^2}{2(1-\varepsilon)K}.
\]
% lean: CausalSmith.PartialID.UnlinkedPropensityAte.worstCaseAmbiguity_finiteLabel_lower

For this same release, \cref{obj:thm:all-label-ambiguity} supplies two causal laws with a common released law and treatment effects separated by \(D_H(g)\). Under their common \(n\)-row law, honesty makes the two coverage events intersect with probability at least \(1-2\alpha\). On that intersection, the interval has length at least \(D_H(g)\). Hence every honest \((g,C_n)\) satisfies
\[
\sup_{P\in\mathfrak P(H,g)}
\mathbb E_{P_{\mathrm{rel}}^{\otimes n}}|C_n|
\geq(1-2\alpha)D_H(g)
\geq A_{\varepsilon,m_f,\alpha}K^{-1},
\qquad
A_{\varepsilon,m_f,\alpha}
:=\frac{(1-2\alpha)m_f\ell^2}{2(1-\varepsilon)}>0.
\]
Taking the infimum over honest pairs preserves this bound. % lean: CausalSmith.PartialID.UnlinkedPropensityAte.minimaxHonestLength_finiteLabel_lower

Define
\[
c_{\varepsilon,m_f,\alpha}
=\frac12\min\{A_{\varepsilon,m_f,\alpha},B_\alpha\}>0.
\]
The label and sampling bounds give
\[
c_{\varepsilon,m_f,\alpha}
\bigl(K^{-1}+n^{-1/2}\bigr)
\leq
\max\{A_{\varepsilon,m_f,\alpha}K^{-1},
      B_\alpha n^{-1/2}\}
\leq\mathcal L^\star_{n,K,\alpha}(H).
\]
% lean: CausalSmith.PartialID.UnlinkedPropensityAte.minimaxHonestLength_rate_assembly

For the upper bound, use the equal-width release and midpoints in the statement. Its endpoint convention assigns every \(e\in\mathcal E\) a label and gives
\[
z_{g_K(e)}\in[\varepsilon,1-\varepsilon],
\qquad |e-z_{g_K(e)}|\leq\frac{\ell}{2K}.
\]
Fix \(P\in\mathfrak P(H,g_K)\), and choose conditional means
\[
\eta_a(e)=\mathbb E_P[Y_a\mid E=e]\in[0,1],
\qquad a\in\{0,1\}.
\]
The score-based assignment and consistency conditions in \cref{obj:def:causal-law-class} give
\[
\begin{aligned}
\mathbb E_{P_{\mathrm{rel}}^{\otimes n}}\widehat T_n
&=\int_{\mathcal E}
 \left\{\frac{e\eta_1(e)}{z_{g_K(e)}}
 -\frac{(1-e)\eta_0(e)}{1-z_{g_K(e)}}\right\}H(de),\\
\theta(P)&=\int_{\mathcal E}\{\eta_1(e)-\eta_0(e)\}\,H(de).
\end{aligned}
\]
Subtracting and using \(z(1-z)\geq\varepsilon(1-\varepsilon)\) on the score interval gives
\[
\begin{aligned}
\mathbb E_{P_{\mathrm{rel}}^{\otimes n}}\widehat T_n-\theta(P)
&=\int_{\mathcal E}(e-z_{g_K(e)})
 \left\{\frac{\eta_1(e)}{z_{g_K(e)}}
       +\frac{\eta_0(e)}{1-z_{g_K(e)}}\right\}H(de),\\
\left|\mathbb E_{P_{\mathrm{rel}}^{\otimes n}}\widehat T_n-\theta(P)\right|
&\leq\frac{\ell}{2K\varepsilon(1-\varepsilon)}.
\end{aligned}
\]
% lean: CausalSmith.PartialID.UnlinkedPropensityAte.midpointWeightedRow_integral_sub_ate_abs_le

Each summand of \(\widehat T_n\) lies in \([-\varepsilon^{-1},\varepsilon^{-1}]\). Hoeffding’s inequality for the independent released rows, with
\[
t_n=\frac4\varepsilon\sqrt{\frac{\log(4/\alpha)}n},
\]
therefore gives
\[
P_{\mathrm{rel}}^{\otimes n}
 \left(\left|\widehat T_n-
 \mathbb E_{P_{\mathrm{rel}}^{\otimes n}}\widehat T_n\right|\geq t_n\right)
\leq 2\exp\!\left(-\frac{n\varepsilon^2t_n^2}{2}\right)
=2(\alpha/4)^8\leq\alpha.
\]
% lean: CausalSmith.PartialID.UnlinkedPropensityAte.midpointWeightedCenter_deviation_probability_le

The radius in the statement is \(R_{n,K}=t_n+\ell/[2K\varepsilon(1-\varepsilon)]\). Outside the displayed deviation event, the bias bound implies
\(\lvert\widehat T_n-\theta(P)\rvert\leq R_{n,K}\). Also \(\theta(P)\in[-1,1]\), since \(Y_0,Y_1\in[0,1]\). Clamping the two endpoints to \([-1,1]\) therefore preserves containment of \(\theta(P)\). The radius is nonnegative, and the clamp is measurable and nondecreasing, so the stated endpoints form an interval for every sample. The release is measurable. As the argument holds for every \(P\in\mathfrak P(H,g_K)\), \((g_K,C_n)\in\mathcal J_{n,K,\alpha}(H)\) by \cref{obj:def:honest-procedures}. % lean: CausalSmith.PartialID.UnlinkedPropensityAte.equalWidth_midpointInterval_honest

Finally, the clamp is one-Lipschitz. Pointwise in the sample,
\[
|C_n|
\leq 2R_{n,K}
=\frac{\ell}{\varepsilon(1-\varepsilon)}K^{-1}
 +\frac8\varepsilon\sqrt{\log(4/\alpha)}\,n^{-1/2}
\leq U_{\varepsilon,\alpha}
 \bigl(K^{-1}+n^{-1/2}\bigr).
\]
Integrating under every released \(n\)-row law gives the asserted expected-length bound with \(C_{\varepsilon,\alpha}\geq U_{\varepsilon,\alpha}\). Taking the supremum over causal laws and the infimum over honest pairs also gives
\[
\mathcal L^\star_{n,K,\alpha}(H)
\leq C_{\varepsilon,\alpha}\bigl(K^{-1}+n^{-1/2}\bigr).
\]
Together with the two lower bounds, this proves the claims. % lean: CausalSmith.PartialID.UnlinkedPropensityAte.minimaxHonestLength_order
\end{proof}

\begin{proof}[Proof of \cref{obj:thm:external-score-root-order}]
\emph{Step 1: Lower certificates.}
\Cref{obj:lem:external-two-source-lower-certificate} gives \(0<S_g\leq1\), the trial and score-log lower bounds, and the stated lower bound on the normalized risk. It also gives \(c_{\mathrm{log}}(g,\alpha)>0\). We construct an interval attaining the upper bound. % lean: CausalSmith.PartialID.UnlinkedPropensityAte.externalExcessRisk_lower_certificates

\emph{Step 2: Empirical distances and endpoint stability.}
Fix positive \(n,m\) and an admissible pair \((P,H)\). For each arm \(a\) and label \(r\), define the population finite measures
\[
\lambda_{ar}(B)=P(R=r,A=a,Y\in B),\quad B\subseteq[0,1],
\qquad
\sigma_{ar}(B)=\int_{B\cap C_r}p_a(e)\,H(de),\quad B\subseteq\mathcal E.
\]
Their total masses agree by the score-marginal, assignment, and release conditions. For finite measures \(\xi,\xi'\) on a compact interval \(\mathcal I\), write
\[
d_{\mathcal I}(\xi,\xi')
=
|\xi(\mathcal I)-\xi'(\mathcal I)|
+\int_{\mathcal I}|F_\xi(t)-F_{\xi'}(t)|\,dt,
\qquad F_\xi(t)=\xi((-\infty,t]).
\]
For a joint sample \(\omega\), set
\[
D_Y(\omega)=\sum_{a,r}d_{[0,1]}(\widehat\lambda_{ar}(\omega),\lambda_{ar}),
\qquad
D_E(\omega)=\sum_{a,r}d_{\mathcal E}(\widehat\sigma_{ar}(\omega),\sigma_{ar}).
\]
The empirical measures are those of \cref{obj:def:external-projection-handle}. Each empirical cumulative mass, including its total mass, is an average of independent variables in \([0,1]\). Its expected absolute deviation is at most \(1/(2\sqrt n)\) for a trial cell and \(1/(2\sqrt m)\) for a score-log cell. Integrating over intervals of lengths \(1\) and \(1-2\varepsilon\), respectively, and summing the \(2J\) cells gives
\[
\mathbb E_{\mathbb Q_{P,H}^{n,m}}D_Y\leq\frac{2J}{\sqrt n},
\qquad
\mathbb E_{\mathbb Q_{P,H}^{n,m}}D_E
\leq\frac{J(2-2\varepsilon)}{\sqrt m}. \label{eq:external-root-order-mean}
\]

Here is the stability estimate used below. For compatible finite measures \(\xi\) on \([0,1]\) and \(\zeta\) on \(\mathcal E\), let their common mass be \(\kappa\). When \(\kappa>0\), define
\[
\Psi_a^-(\xi,\zeta)
=\kappa\int_0^1 Q_{\xi/\kappa}(u)
 Q_{(p_a^{-1})_\#(\zeta/\kappa)}(1-u)\,du,
\qquad
\Psi_a^+(\xi,\zeta)
=\kappa\int_0^1 Q_{\xi/\kappa}(u)
 Q_{(p_a^{-1})_\#(\zeta/\kappa)}(u)\,du;
\]
when \(\kappa=0\), set both values to zero. For another compatible pair \((\xi',\zeta')\) of mass \(\kappa'\), put \(\kappa_\wedge=\min\{\kappa,\kappa'\}\). Suppose first that both masses are positive. The quantile transport identity on a compact interval \(\mathcal I\) identifies the integral of the absolute difference of normalized quantiles with the integral of the absolute difference of normalized CDFs. If \(\kappa\leq\kappa'\), then \(0\leq F_{\xi'}(t)\leq\kappa'\) gives, pointwise,
\[
\kappa\left|\frac{F_\xi(t)}{\kappa}-\frac{F_{\xi'}(t)}{\kappa'}\right|
\leq |F_\xi(t)-F_{\xi'}(t)|+|\kappa-\kappa'|.
\]
The same inequality with the pairs exchanged covers \(\kappa'\leq\kappa\). Applying it to the outcome and score CDFs and integrating over supports of lengths \(1\) and \(1-2\varepsilon\), respectively, yields
\[
\begin{aligned}
\kappa_\wedge\int_0^1
 |Q_{\xi/\kappa}(u)-Q_{\xi'/\kappa'}(u)|\,du
&\leq
\int_0^1|F_\xi(t)-F_{\xi'}(t)|\,dt+|\kappa-\kappa'|,\\
\kappa_\wedge\int_0^1
 |Q_{\zeta/\kappa}(u)-Q_{\zeta'/\kappa'}(u)|\,du
&\leq
\int_{\mathcal E}|F_\zeta(t)-F_{\zeta'}(t)|\,dt
 +(1-2\varepsilon)|\kappa-\kappa'|.
\end{aligned} \label{eq:external-root-order-transport}
\]
Since \(0\leq y\leq1\), \(p_a(e)\geq\varepsilon\), and \(e\mapsto p_a(e)^{-1}\) is \(\varepsilon^{-2}\)-Lipschitz, the product difference in either definition of \(\Psi_a^\pm\), multiplied by \(\kappa_\wedge\), is bounded using \cref{eq:external-root-order-transport} by
\[
\varepsilon^{-1}
 \left(\int_0^1|F_\xi-F_{\xi'}|+|\kappa-\kappa'|\right)
+\varepsilon^{-2}
 \left(\int_{\mathcal E}|F_\zeta-F_{\zeta'}|
       +|\kappa-\kappa'|\right). \label{eq:external-root-order-product}
\]
The remaining mass contributes at most \(\varepsilon^{-1}|\kappa-\kappa'|\). Because both distances contain that mass difference, \cref{eq:external-root-order-product} implies, with
\(\Lambda_\varepsilon=\varepsilon^{-1}+\varepsilon^{-2}\),
\[
|\Psi_a^\pm(\xi,\zeta)-\Psi_a^\pm(\xi',\zeta')|
\leq
\Lambda_\varepsilon
 \bigl(d_{[0,1]}(\xi,\xi')+d_{\mathcal E}(\zeta,\zeta')\bigr). \label{eq:external-root-order-stability}
\]
If either mass is zero, the bound \(0\leq\Psi_a^\pm(\xi,\zeta)\leq\varepsilon^{-1}\kappa\) and the mass term in the distance give \cref{eq:external-root-order-stability} directly. Summing \cref{eq:external-root-order-stability} over labels gives the same bound, with summed distances, for each arm endpoint. By \cref{obj:lem:fixed-released-law-baseline}, those endpoints at \((\lambda_{ar},\sigma_{ar})_{a,r}\) are \(\underline\mu_a,\overline\mu_a\). % lean: CausalSmith.PartialID.UnlinkedPropensityAte.externalJointTrial_meanDistance_le, CausalSmith.PartialID.UnlinkedPropensityAte.externalJointLog_meanDistance_le, CausalSmith.PartialID.UnlinkedPropensityAte.projectedCellSummand_abs_le, CausalSmith.PartialID.UnlinkedPropensityAte.projectedArmEndpoint_population_eq

\emph{Step 3: Honesty and excess length.}
Use the interval \(C_{n,m}\) of \cref{obj:def:external-projection-handle}. Its countable rational sieve and selected endpoints give measurable endpoints and a closed interval in \([-1,1]\). The release \(g:\mathcal E\to\mathcal R_J\), with nonempty \(\mathcal E\), also ensures \(J\geq1\). The radii in that construction are
\[
r_n=\frac{4J}{\alpha\sqrt n},
\qquad
r_m=\frac{2J(2-2\varepsilon)}{\alpha\sqrt m}.
\]
By \cref{eq:external-root-order-mean}, Markov's inequality applied to the nonnegative distances gives
\[
\mathbb Q_{P,H}^{n,m}\{D_Y<r_n,\ D_E<r_m\}\geq1-\alpha. \label{eq:external-root-order-coverage}
\]
On the strict-ball event in \cref{eq:external-root-order-coverage}, define the available slack
\[
\eta_* = \min\{r_n-D_Y(\omega),\,r_m-D_E(\omega)\}>0,
\qquad \eta_k=\frac{\eta_*}{k+1}\quad(k\geq1).
\]
For each arm-label cell, its population outcome and score measures have the same finite mass. Given \(\eta>0\), choose a common nonnegative rational mass and a common positive atom count for two equal-weight atomic approximations, with outcome and score atoms on their respective dense rational grids. To see the quantitative bound, first approximate each normalized measure by equally weighted quantile atoms; increasing the common atom count makes the sum of their integrated CDF errors arbitrarily small. Moving those atoms to the rational grids preserves that convergence. If the common population mass is \(\kappa\) and the chosen rational mass is \(\kappa'\), changing the weight of each atom from \(\kappa/N\) to \(\kappa'/N\) adds at most \((1+|\mathcal I|)|\kappa-\kappa'|\) to the distance on support \(\mathcal I\). Thus the combined weight error for the two supports is at most \((4-2\varepsilon)|\kappa-\kappa'|\). Choosing \(\kappa'\) sufficiently close to \(\kappa\) gives combined cell error below any prescribed tolerance; this also covers zero mass by taking rational mass zero.

Apply this construction in every cell with tolerance \(\eta_k/[2(J+1)]\). Since there are \(2J\) cells, it gives a compatible rational array \(b_k\) satisfying
\[
\sum_{a,r}\!\left(d_{[0,1]}(\lambda_{ar},\lambda_{ar}(b_k))
+d_{\mathcal E}(\sigma_{ar},\sigma_{ar}(b_k))\right)
<\frac{2J}{2(J+1)}\eta_k<\eta_k.
\]
Here \(\lambda_{ar}(b_k)\) and \(\sigma_{ar}(b_k)\) denote the two measures encoded by \(b_k\). Each of the two separate distance sums is therefore below \(\eta_k\). The triangle inequality gives empirical distances below \(D_Y(\omega)+\eta_k<r_n\) and \(D_E(\omega)+\eta_k<r_m\), so \(b_k\) belongs to both closed empirical balls. As \(k\to\infty\), \cref{eq:external-root-order-stability} makes all four arm endpoints converge to their population values. Write
\[
I(b_k)=[\underline\theta_k,\overline\theta_k],
\qquad
I(P_{\mathrm{rel}};H,g)=[\underline\theta_\star,\overline\theta_\star].
\]
Both candidate endpoints converge to their corresponding population endpoints. Fix any \(v\in I(P_{\mathrm{rel}};H,g)\) and clamp it to the candidate interval:
\[
v_k=\min\{\max\{v,\underline\theta_k\},\overline\theta_k\}\in I(b_k).
\]
If \(v<\underline\theta_k\), the distance \(|v_k-v|\) is at most \(|\underline\theta_k-\underline\theta_\star|\); if \(\overline\theta_k<v\), it is at most \(|\overline\theta_k-\overline\theta_\star|\); otherwise it is zero. Thus \(v_k\to v\). Every point of the population interval therefore lies in the closure of the union of candidate intervals, and hence in their closed convex hull. Write \(\theta(P)=\mathbb E_P[Y_1-Y_0]\). By \cref{obj:lem:fixed-released-law-baseline}, \(\theta(P)\) belongs to that sharp interval; its potential-outcome support also places it in \([-1,1]\). Thus the clipped hull contains \(\theta(P)\) on the event in \cref{eq:external-root-order-coverage}, proving
\[
C_{n,m}\in\mathcal J^{\mathrm{ext}}_{n,m,\alpha}(g),
\qquad
C_{n,m}(\omega)=\text{the closed interval of \cref{obj:def:external-projection-handle} at }\omega
\quad\text{for every }\omega. \label{eq:external-root-order-honesty}
\]

For any retained candidate, the two empirical-ball inequalities and the triangle inequality bound its summed distance from the population array by \(r_n+D_Y(\omega)+r_m+D_E(\omega)\). \Cref{eq:external-root-order-stability} therefore bounds each candidate arm endpoint's error by
\(\Lambda_\varepsilon(r_n+D_Y(\omega)+r_m+D_E(\omega))\). Each average-treatment-effect endpoint combines two arm endpoints; taking the hull and clipping it to \([-1,1]\) yields
\[
\left(\operatorname{length}(C_{n,m}(\omega))
-\operatorname{length}I(P_{\mathrm{rel}};H,g)\right)_+
\leq
4\Lambda_\varepsilon
 \bigl(r_n+D_Y(\omega)+r_m+D_E(\omega)\bigr). \label{eq:external-root-order-excess}
\]
In either fallback branch of the construction the interval is \(\{0\}\), so its positive excess is zero and \cref{eq:external-root-order-excess} holds there as well. Taking expectations and using \cref{eq:external-root-order-mean} gives, uniformly over admissible \((P,H)\),
\[
\mathbb E_{\mathbb Q_{P,H}^{n,m}}
 \left[\left(\operatorname{length}(C_{n,m})
-\operatorname{length}I(P_{\mathrm{rel}};H,g)\right)_+\right]
\leq
4\Lambda_\varepsilon
 \left(
 \frac{4J/\alpha+2J}{\sqrt n}
 +\frac{2J(2-2\varepsilon)/\alpha+J(2-2\varepsilon)}{\sqrt m}
 \right). \label{eq:external-root-order-mean-excess}
\]
The joint marginal and independence assumptions identify \(\mathbb Q_{P,H}^{n,m}\) with the product experiment in \cref{obj:def:external-experiment}. Hence honesty and \cref{eq:external-root-order-mean-excess}, followed by the infimum and supremum in \cref{obj:def:external-excess-risk}, give the same upper bound for \(\mathcal R^\star_{n,m,\alpha}(g)\). Set
\[
C_{\varepsilon,J,\alpha,g}
=
4\Lambda_\varepsilon
 \left(\frac{4J}{\alpha}+2J
       +\frac{2J(2-2\varepsilon)}{\alpha}
       +J(2-2\varepsilon)\right)+1>0.
\]
Both \cref{eq:external-root-order-mean-excess} and the minimax bound are then at most
\(C_{\varepsilon,J,\alpha,g}(n^{-1/2}+m^{-1/2})\), as required. % lean: CausalSmith.PartialID.UnlinkedPropensityAte.externalProjectionIntervalProcedure_honest, CausalSmith.PartialID.UnlinkedPropensityAte.externalProjectionIntervalProcedure_Icc, CausalSmith.PartialID.UnlinkedPropensityAte.externalProjectionIntervalProcedure_population_excessLength_le, CausalSmith.PartialID.UnlinkedPropensityAte.externalProjectionIntervalProcedure_population_meanExcessLength_le, CausalSmith.PartialID.UnlinkedPropensityAte.externalExcessRisk_le_of_honest_uniform_bound

\emph{Step 4: Normalized bounds.}
For positive \(n,m\), the upper bound just proved and \cref{obj:def:external-normalized-risk} imply
\[
0\leq T_{n,m,\alpha}(g)
=b_{n,m}\mathcal R^\star_{n,m,\alpha}(g)
\leq C_{\varepsilon,J,\alpha,g},
\qquad
b_{n,m}=(n^{-1/2}+m^{-1/2})^{-1}. \label{eq:external-root-order-normalized}
\]
The lower inequality follows from the positive trial certificate in \cref{obj:lem:external-two-source-lower-certificate}. Along the nonempty joint regime \(n,m\to\infty\), \(n/m\to\rho>0\), \cref{eq:external-root-order-normalized} bounds the limsup by \(C_{\varepsilon,J,\alpha,g}\) and ensures that the liminf is at most the limsup. The same cited lower certificate supplies
\[
\liminf T_{n,m,\alpha}(g)\geq
\max\left\{
\frac{c_{\mathrm{tr}}(\alpha)}{1+\sqrt\rho},
\frac{c_{\mathrm{log}}(g,\alpha)\sqrt\rho}{1+\sqrt\rho}
\right\}>0.
\]
Together with the score-log liminf bound in \cref{obj:lem:external-two-source-lower-certificate}, these are all the asserted inequalities. % lean: CausalSmith.PartialID.UnlinkedPropensityAte.externalExcessRisk_root_order
\end{proof}

\bibliographystyle{plainnat}

\bibliography{references}

\end{document}
